\documentclass[../main.tex]{subfiles}
\graphicspath{{\subfix{../}}}
\begin{document}

Hullámegyenletek származtatása, megoldásai, rugalmas és elektromágneses hullámok. Diszperzió, csoport- és fázissebesség, Doppler-effektus.
Interferencia, diffrakció. Elektromágneses hullámok terjedése vákuumban, dielektrikumban és vezetőkben. Polarizáció. Retardált potenciálok. Dipólsugárzás.

\subsection{Hullámegyenletek származtatása, megoldásai}
\subsubsection{A hullámegyenlet}
Mi a "hullám"? A hullám egy nehezen megfogható jelenség, mert a koncepció alapvetően kicsit ködös. De általánosságban elmondható, hogy a hullám egy zavar egy folytonos közegben,
ami képes terjedni egy fix alakban és konstans sebességgel. Ez a definíció nem fed le minden esetet, de kiindulási alapnak megfelelő.

De hogyan lehet ezt az objektumot matematikailag reprezentálni? A "zavart" magát jellemezhetjük egy skalár fügvénnyel ($f(z, t)$), ami a közeg minden pontjához hozzárendel egy értéket.
Erről a függvényről tudjuk, hogy $v$ sebességgel terjed a $z$ tengely mentén, de mivel az alakja állandó, ezért $t=0$-ban (itt legyen a függvény $f(0, t) = g(z)$) a függvény értéke ugyanaz lesz mint $t$ időpillanatban a $z - vt$ helyen:
\begin{equation}
    f(z, t) = f(z-vt, 0) = g(z - vt)
\end{equation}
Ez az állítás jellemzi matematikailag a hullámok lényegét.  Ez elmondja, hogy a $f$ függvény, bár függhet $z$-től és $t$-től is, de csak egy nagyon speciális kombinációban: $z - vt$.
Tehát csak olyan függvényeket tekinthetünk hullámnak, amiben teljesül ez a feltétel. Tehát ezek a függvények lehetnek hullámok:
\begin{equation}
    f(z, t) = A \sin(k(z - vt)) \qquad f(z, t) = \frac{1}{b(z-vt)^2 + 1} \qquad f(z, t) = A e^{i k (z - vt)}
\end{equation}
De ezek a függvények NEM lehetnek hullámok:
\begin{equation}
    f(z, t) = Ae^{-b(bz^2 + vt)} \qquad f(z, t) = A \sin(bz) \cos(bvt)^3
\end{equation}

\subsubsection{Klasszikus hullám}
A klasszikus mechanikus hullámok esetében valóban van egy közeg amiben a hullám terjed. Vegyünk például egy kötelet ami a $z$ tengely mentén van kifeszítve, és az egyik végét $t=0$ időpillanatban megmozgatjuk.
\begin{figure}[H]
    \centering
    \includegraphics[width=0.5\textwidth]{imgs/8-tetel/kotel_hullam.png}
    \caption{Kötél hulláma}
\end{figure}
Tegyük fel, hogy a kötél $T$ feszültség alatt van, és $\mu$ lineáris tömegsűrűséggel rendelkezik. Ekkor meghatározhatjuk az erőt ami hat egy kis $dz$ hosszúságú darabra ($z$ és $dz$ pontok között) a kötélen hat, amikor az egyensúlyból kicsit kimozdítjuk:
\begin{equation}
    \Delta F = T \sin \theta' - T \sin \theta
\end{equation}
Ahol $\theta'$ az a szög ami a $z + dz$ pontban lévő kötéldarabra jellemző, és $\theta$ a $z$ pontban lévő kötéldarabra jellemző. Ha a kimozdulás kicsi, akkor a szögek is kicsik lesznek, és akkor a szinuszokat kicserélhetjük tangensre:
\begin{equation}
\Delta F \approx T \tan \theta' - T \tan \theta = T \frac{\partial f}{\partial z} \Big|_{z+dz} - T \frac{\partial f}{\partial z} \Big|_{z} = T \frac{\partial^2 f}{\partial z^2} dz
\end{equation}
De ugyanezt az erőt a Newton II. törvénye alapján is meghatározhatjuk. Ha a kötélnek egy egységen $\mu$ a tömege, akkor az erő a következőképpen írható fel a $dz$ hosszúságú kötéldarabra:
\begin{equation}
   \Delta F = \mu dz \frac{\partial^2 f}{\partial t^2}
\end{equation}
Tehát a két egyenletnek egyenlőnek kell lennie:
\begin{equation}
    T \frac{\partial^2 f}{\partial z^2} dz = \mu dz \frac{\partial^2 f}{\partial t^2}
\end{equation}
Ahol a $dz$-k kiesnek, és a következőt kapjuk:
\begin{equation}
    \frac{\partial^2 f}{\partial z^2} = \frac{\mu}{T} \frac{\partial^2 f}{\partial t^2}
\end{equation}
Ez a \textbf{klasszikus hullámegyenlet} egy dimenzióban. Ha bevezetjük a hullámsebességet:
\begin{equation}
    v = \sqrt{\frac{T}{\mu}}
\end{equation}
Akkor a hullámegyenlet a következő formát ölti:
\begin{equation}
    \boxed{\frac{\partial^2 f}{\partial z^2} = \frac{1}{v^2} \frac{\partial^2 f}{\partial t^2}}
\end{equation}
Ez az egyenlet írja le a klasszikus mechanikai hullámokat egy dimenzióban. Az egyenlet megoldásai azok a függvények, amik kielégítik a hullám feltételt, azaz $f(z, t) = g(z - vt)$ formában vannak.
Bizonyítás:
\begin{equation}
    \frac{\partial f}{\partial z} = \frac{dg}{du} \frac{\partial u}{\partial z} = \frac{dg}{du}, \qquad \frac{\partial f}{\partial t} = \frac{dg}{du} \frac{\partial u}{\partial t} = -v \frac{dg}{du}
\end{equation}
A második deriváltak pedig:
\begin{equation}
    \begin{aligned}
        \frac{\partial^2 f}{\partial z^2} &= \frac{d}{dz} \left( \frac{dg}{du} \right) = \frac{d^2 g}{du^2} \frac{\partial u}{\partial z} = \frac{d^2 g}{du^2} \\
        \frac{\partial^2 f}{\partial t^2} &= \frac{d}{dt} \left( -v \frac{dg}{du} \right) = -v \frac{d^2 g}{du^2} \frac{\partial u}{\partial t} = v^2 \frac{d^2 g}{du^2}
    \end{aligned}
\end{equation}
Behelyettesítve a hullámegyenletbe:
\begin{equation}
    \frac{\partial^2 f}{\partial z^2} = \frac{1}{v^2} \frac{\partial^2 f}{\partial t^2} \qquad \Rightarrow \qquad \frac{d^2 g}{du^2} = \frac{1}{v^2} \left( v^2 \frac{d^2 g}{du^2} \right)
\end{equation}
Tehát a hullám feltétel kielégíti a hullámegyenletet. $g(u)$ bármilyen differenciálható függvény lehet. Felmerülhet a kérdés, hogy hogyan lehetséges ez, hiszen egy másodrendű differenciálegyenletnek
két független megoldása kell legyen. A válasz az, hogy ez a megkötés csak sima differenciáltakra igaz, parciális deriváltaknak nem kell ennek a megkötésnek megfelelnie! De nem csak ezek oldják meg
az egyenletet, hanem a $g(z + vt)$ formájú függvények is, amik a negatív irányba terjedő hullámokat írják le. Tehát a hullámegyenlet általános megoldása a következőképpen írható fel:
\begin{equation}
    \boxed{f(z, t) = g(z - vt) + h(z + vt)}
\end{equation}
Ahol $g$ és $h$ tetszőleges differenciálható függvények. Ebből az is jól látszik, hogy a hullámegyenlet lineáris, hiszen két megoldás összege is megoldás.

A leggyakabban használt megoldás a szinuszos hullám, ami a következőképpen írható fel:
\begin{equation}
    f(z, t) = A \sin(kz - \omega t + \phi)
\end{equation}
Ahol $A$ a hullám amplitúdója, $k$ a hullámszám, $\omega$ a körfrekvencia, és $\phi$ a kezdeti fázis. Ezek a mennyiségek a következőképpen kapcsolódnak egymáshoz:
\begin{equation}
    k = \frac{2 \pi}{\lambda} \qquad \omega = 2 \pi f = kv \qquad v = \frac{\omega}{k} = \lambda f
\end{equation}
Ahol $\lambda$ a hullámhossz, $k$ a hullámszám, és $f$ a frekvencia.

A másik gyakan használt megoldás a komplex hullám:
\begin{equation}
    f(z, t) = A e^{i(kz - \omega t)}
\end{equation}
Ez a megoldás is kielégíti a hullámegyenletet, és a valós részét véve visszakapjuk a szinuszos hullámot (ha $\phi = 0$). A komplex exponenciális forma kényelmesebb a számolások során.

Nagyon fontos tulajdonsága a szinuszos hullámoknak, hogy bármelyik hullám kifejezhető szinuszos hullámok lineáris kombinációjaként! Tehát egy tetszőleges hullámfüggvényt fel tudjuk írni a következő alakban:
\begin{equation}
    \bar{f}(z, t) = \int_{-\infty}^{\infty} \bar{A}(k) e^{i(kz - \omega t)} \, dk
\end{equation}
Ez a Fourier-transzformáció, ami segítségével bármilyen hullámot fel tudunk bontani komponenseire.

\subsubsection{Határfeltételek: Visszaverődés és transzmisszió}
Eddig azt feltételeztük, hogy a hullám egy végtelen közegben terjed. De mi történik akkor, ha a hullám egy határfelülethez érkezik, ahol a közeg megváltozik?
Például tegyük fel, hogy a kötelet egy másik kötélhez kötöttük, aminek azonos a feszültsége, de eltérő a lineáris tömegsűrűsége. Ekkor a hullám terjedési sebessége más lesza két kötélben (hiszen $v = \sqrt{T/\mu}$).
Az egyszerűség kedvéért tegyük a közési pontot a $z=0$ helyre. Ekkor a hullám a $z<0$ tartományból érkezik, és a következőképpen írható fel:
\begin{equation}
    \bar{f}(z, t) = \bar{A}_1 e^{i(k_1 z - \omega t)} \qquad z < 0
\end{equation}
Ekkor a bal oldalól megjelenik egy \textbf{visszavert hullám} is, ami a következőképpen írható fel:
\begin{equation}
    \bar{f}_r(z, t) = \bar{A}_R e^{i(-k_1 z - \omega t)} \qquad z < 0
\end{equation}
A jobb oldalra pedig egy \textbf{átvitt hullám} is megjelenik, ami a következőképpen írható fel:
\begin{equation}
    \bar{f}_t(z, t) = \bar{A}_T e^{i(k_2 z - \omega t)} \qquad z > 0
\end{equation}
Ami folytatódik a második kötélben. Itt $k_1$ és $k_2$ a két kötél hullámszámai, amik a következőképpen kapcsolódnak a frekvenciához:
\begin{equation}
    k_1 = \frac{\omega}{v_1} = \omega \sqrt{\frac{\mu_1}{T}} \qquad k_2 = \frac{\omega}{v_2} = \omega \sqrt{\frac{\mu_2}{T}}
\end{equation}
Mivel a hullám sebesség eltér a két kötélen, ezért a hullámszám és a hullámhossz is eltérő lesz:
\begin{equation}
    \frac{\lambda_1}{\lambda_2} = \frac{k_2}{k_1} = \frac{v_1}{v_2}
\end{equation}
Tehát egy szinuszos hullám egy határfelületnél:
\begin{equation}
    \bar{f}(z, t) = \begin{cases}
        \bar{A}_1 e^{i(k_1 z - \omega t)} + \bar{A}_R e^{i(-k_1 z - \omega t)} & z < 0 \\
        \bar{A}_T e^{i(k_2 z - \omega t)} & z > 0
    \end{cases}
\end{equation}
Ha a kapcsolódási pont $z=0$ helyen van. A zavarnak egy kicsit balra ($z = 0^-$) ugyanannyinak kell lennie mint egy kicsit jobbra ($z = 0^+$), hiszen a kötél folytonos:
\begin{equation}
    \bar{f}(0^-, t) = \bar{f}(0^+, t)
\end{equation}
Ha a köteleknek elhanyagolható a tömege, akkor a derviáltjuk is ugyanaz:
\begin{equation}
    \frac{\partial \bar{f}}{\partial z} \Big|_{0^-} = \frac{\partial \bar{f}}{\partial z} \Big|_{0^+}
\end{equation}
különben lenne egy nettó erő ami erre a kötéldarabra hat és végtelenül gyorsulna. Ezek a határfeltételek a valós hullámfüggvényekre igazak, de mivel a komplex exponenciális imaginárius része csak egy $\pi/2$ fáziseltolódásban tér el,
ezért a komplex hullámokra is ezek lesznek a határfeltételek.A határfeltételekből az amplitúdókra a következő összefüggések adódnak:
\begin{equation}
    \begin{aligned}
        \bar{A}_1 + \bar{A}_R &= \bar{A}_T \\
        k_1 (\bar{A}_1 - \bar{A}_R) &= k_2 \bar{A}_T
    \end{aligned}
\end{equation}
Amiből kifejezve az amplitúdókat:
\begin{equation}
    \begin{aligned}
        \bar{A}_R &= \frac{k_1 - k_2}{k_1 + k_2} \bar{A}_1 \\
        \bar{A}_T &= \frac{2 k_1}{k_1 + k_2} \bar{A}_1
    \end{aligned}
\end{equation}
Ezek az amplitúdók arányai a visszavert és az átvitt hullámoknak az eredeti hullám amplitúdójához képest. Ugyanezt a hullám sebességével is kifejezhetjük:
\begin{equation}
    \begin{aligned}
        \bar{A}_R &= \frac{v_2 - v_1}{v_2 + v_1} \bar{A}_1 \\
        \bar{A}_T &= \frac{2 v_2}{v_2 + v_1} \bar{A}_1
    \end{aligned}
\end{equation}
A valós amplitúdók és fázisok pedig:
\begin{equation}
    A_R e^{i \phi_R} = \frac{v_2 - v_1}{v_2 + v_1} A_1 e^{i \phi_1} \qquad A_T e^{i \phi_T} = \frac{2 v_2}{v_2 + v_1} A_1 e^{i \phi_1}
\end{equation}

Tehát ha például a második kötél könnyebb mint az első, akkor a második kötélen nagyobb a hullámsebesség ($v_2 > v_1$), de a fázis azonos lesz minden hullámnál ($\phi_R = \phi_T = \phi_1$).
Tehát az amplitúdók a következőképpen alakulnak:
\begin{equation}
    A_R = \frac{v_2 - v_1}{v_2 + v_1} A_1 \qquad A_T = \frac{2 v_2}{v_2 + v_1} A_1
\end{equation}
Ha viszont a második kötél nehezebb mint az első, akkor a második kötélen kisebb a hullámsebesség ($v_2 < v_1$), és a visszavert hullám fázisa eltolódik $\pi$-val ($\phi_R = \phi_1 + \pi$), míg az átvitt hullám fázisa
azonos marad ($\phi_T = \phi_1$). Az amplitúdók pedig a következőképpen alakulnak:
\begin{equation}
    A_R = \frac{v_1 - v_2}{v_2 + v_1} A_1 \qquad A_T = \frac{2 v_2}{v_2 + v_1} A_1
\end{equation}
Mivel (ha $cos$-al adjuk meg a hullámot):
\begin{equation}
    \cos(- k_1z - \omega t + \phi_1 + \pi) = - \cos(-k_1z - \omega t + \phi_1)
\end{equation}
Egy speciális eset ha a második kötél végtelen nehéz - vagy egy elmozdithatatlan tárgyhoz van a kötél hozzákötve (pl. egy fal), ekkor:
\begin{equation}
    A_R = A_1 \qquad \phi_R = \phi_1 + \pi \qquad A_T = 0
\end{equation}
Tehát a hullám teljesen visszaverődik, és fáziseltolódik $\pi$-val.

\subsection{Polarizáció}
Azokat a hullámokat amiket eddig vizsgáltunk, \textbf{transzverzális hullámoknak} nevezzük, mert a kimozdulás iránya merőleges a hullám terjedési irányára. De léteznek olyan hullámok is, amiknél a kimozdulás iránya párhuzamos a hullám terjedési irányával.
Ezeket a hullámokat \textbf{longitudinális hullámoknak} nevezzük. Ilyen hullámok például a hanghullámok, amik sűrűség és nyomás ingadozásokat hoznak létre a közegben. A hanghullámok terjedése során a levegő részecskéi előre-hátra mozognak
a hullám terjedési irányában, létrehozva a hanghullámot.

De természetesen 3D-ben két irány is lehet merőleges a hullám terjedési irányára. Például egy $z$ tengely mentén terjedő hullám esetében a kimozdulás lehet $x$ vagy $y$ irányú is. Ezeket az irányokat \textbf{polarizációs irányoknak} nevezzük.
A polarizáció azt írja le, hogy az adott hullám melyik irányt preferálja:
\begin{figure}[H]
    \centering
    \includegraphics[width=0.5\textwidth]{imgs/8-tetel/hullam_polarizacio.png}
    \caption{Polarizáció}
\end{figure}
Pl. ha egy kötelet rángatunk, akkor rángathatjuk fel-le, vagy jobbra-balra, vagy akár körkörösen is. Ezek a különböző rángatási irányok különböző polarizációkat hoznak létre a hullámban.
A lineárisan polarizált hullámok esetében a kimozdulás egy adott irányban történik (pl. csak fel-le). A körkörösen polarizált hullámok esetében a kimozdulás iránya folyamatosan változik, és egy kör mentén forog.
Léteznek polarizálatlan hullámok is (pl. a napból érkező fény), ahol nincs egy kitüntetett irány, hanem kvázi random rezeg minden irányba. A polarizációt matematikailag a következőképp írhatjuk le:
\begin{equation}
    \begin{aligned}
        \tilde{\bold{f}}(z, t) &= \tilde{A} e^{i(kz - \omega t)}\hat{\bold{x}} \qquad \text{vertikális polarizáció}\\
        \tilde{\bold{f}}(z, t) &= \tilde{A} e^{i(kz - \omega t)}\hat{\bold{y}} \qquad \text{horizontális polarizáció}\\
        \tilde{\bold{f}}(z, t) &= \tilde{A} e^{i(kz - \omega t)} \hat{\bold{n}} \qquad \text{általános lineáris polarizáció}\\
        \tilde{\bold{f}}(z, t) &= \tilde{A} \left( e^{i(kz - \omega t)}\hat{\bold{x}} + i e^{i(kz - \omega t)}\hat{\bold{y}} \right) \qquad \text{körkörös polarizáció}
    \end{aligned}
\end{equation}
Ahol $\hat{\bold{n}}$ a \textbf{polarizációs vektor}, ami megadja a kimozdulás irányát a $x$ és $y$ tengelyekhez képest. Ez a vektor mindig merőleges a hullám terjedési irányára:
\begin{equation}
    \hat{\bold{n}} \cdot \hat{\bold{z}} = 0
\end{equation}
Azt a szöget amit a polarizációs vektor bezár az $x$ tengellyel, a \textbf{polarizációs szögnek} hívunk. Ha a polarizációs szög $0^\circ$, akkor a hullám vertikálisan polarizált, ha $90^\circ$, akkor horizontálisan polarizált:
\begin{equation}
    \hat{\bold{n}} = \cos \theta_p \hat{\bold{x}} + \sin \theta_p \hat{\bold{y}}
\end{equation}
Ha egy hullám több irányba rezeg, akkor az felgható úgy is, mint több polarizáció szuperpozíciója. Például egy általános lineárisan polarizált hullám felírható úgy is, hogy két különböző irányba polarizált hullám összege:
\begin{equation}
    \tilde{\bold{f}}(z, t) = (\tilde{A} \cos \theta_p) e^{i(kz - \omega t)}\hat{\bold{x}} + (\tilde{A} \sin \theta_p) e^{i(kz - \omega t)}\hat{\bold{y}}
\end{equation}
Ahol $\tilde{A}_x$ és $\tilde{A}_y$ a két irány amplitúdói. Ha ezek az amplitúdók megegyeznek és $\pi/2$ fáziseltolódás van közöttük, akkor körkörös polarizáció jön létre:
\begin{equation}
    \tilde{\bold{f}}(z, t) = \tilde{A} e^{i(kz - \omega t)}\hat{\bold{x}} + \tilde{A} e^{i(kz - \omega t + \frac{\pi}{2})}\hat{\bold{y}}
\end{equation}

\subsection{Elektromágneses hullámok}
\subsubsection{Elektromágneses hullám vákuumban}
Egy olyan térben ahol vákuum van, a Maxwell-egyenletek a következőképpen alakulnak:
\begin{equation}
    \begin{aligned}
        \boldsymbol{\nabla} \cdot \bold{E} &= 0 \qquad \boldsymbol{\nabla} \times \bold{E} = -\frac{\partial \bold{B}}{\partial t} \\
        \boldsymbol{\nabla} \cdot \bold{B} &= 0 \qquad \boldsymbol{\nabla} \times \bold{B} = \mu_0 \epsilon_0 \frac{\partial \bold{E}}{\partial t} 
    \end{aligned}
\end{equation}
Ezek egy kapcsolt, elsőrendű parciális differenciálegyenlet rendszert alkotnak a $\bold{E}$ és $\bold{B}$ terekre. A kapcsoltságot meg lehet szüntetni, ha mindkét egyenletet rotáljuk:
\begin{equation}
    \begin{aligned}
        \boldsymbol{\nabla} \times (\boldsymbol{\nabla} \times \bold{E}) &= \boldsymbol{\nabla} (\boldsymbol{\nabla} \cdot \bold{E}) - \boldsymbol{\nabla}^2 \bold{E} = \boldsymbol{\nabla} \times \left(-\frac{\partial \bold{B}}{\partial t}\right) \\
        &= -\frac{\partial}{\partial t} (\boldsymbol{\nabla} \times \bold{B}) = -\mu_0 \epsilon_0 \frac{\partial^2 \bold{E}}{\partial t^2} \\
        \boldsymbol{\nabla} \times (\boldsymbol{\nabla} \times \bold{B}) &= \boldsymbol{\nabla} (\boldsymbol{\nabla} \cdot \bold{B}) - \boldsymbol{\nabla}^2 \bold{B} = \boldsymbol{\nabla} \times \left( \mu_0 \epsilon_0 \frac{\partial \bold{E}}{\partial t} \right) \\
        &= \mu_0 \epsilon_0 \frac{\partial}{\partial t} (\boldsymbol{\nabla} \times \bold{E}) = -\mu_0 \epsilon_0
    \end{aligned}
\end{equation}
De mivel $\boldsymbol{\nabla} \cdot \bold{E} = 0$ és $\boldsymbol{\nabla} \cdot \bold{B} = 0$, ezért a bal oldali első tagok kiesnek, és a következőket kapjuk:
\begin{equation}
    \boxed{\boldsymbol{\nabla}^2 \bold{E} = \mu_0 \epsilon_0 \frac{\partial^2 \bold{E}}{\partial t^2}} \qquad \boxed{\boldsymbol{\nabla}^2 \bold{B} = \mu_0 \epsilon_0 \frac{\partial^2 \bold{B}}{\partial t^2}}
\end{equation}
Ezek az egyenletek a \textbf{hullámegyenletek} az elektromágneses térre vákuumban. Ezek azt jelentik, hogy az elektromágneses térnek vannak hullámszerű megoldásai, amik a Maxwell-egyenleteknek is megfelelnek.
Ezzel tehát megszüntettük a $\bold{E}$ és $\bold{B}$ terek közötti kapcsoltságát, viszont cserébe másodrendű differenciálegyenleteket kaptunk. Ezek az egyenletek leírják az elektromágneses hullámok terjedését vákuumban,
mivel kielégítik a háromdimenziós hullámegyenletet:
\begin{equation}
    \boldsymbol{\nabla}^2 f = \frac{1}{v^2} \frac{\partial^2 f}{\partial t^2}
\end{equation}
Ahol $v$ a hullám terjedési sebessége. Az elektromágneses hullámok esetében a hullámsebesség a következőképpen adódik:
\begin{equation}
    v = \frac{1}{\sqrt{\mu_0 \epsilon_0}} = c = 3 \times 10^8 \, \text{m/s}
\end{equation}
Tehát az elektromágneses hullámok vákuumban a fénysebességgel terjednek.

\subsubsection{Monokromatikus hullámok}
Az elektromágneses hullámok megoldásai hasonlóak a korábban tárgyalt mechanikai hullámokéhoz, ezért itt is szinuszos hullámalakokat keresünk. Ezeket a hullámokat monokromatikus hullámoknak nevezzük,
mert az elektromágneses hullámok színe a frekvenciától függ, ezét ha egy frekvenciájú hullámot nézünk, akkor az monokromatikus hullám lesz. Ha pedig nincs $x$ és $y$ függése, akkor \textbf{planáris hullámnak} bevezzük. A számunkra érdekes alakja az elektromos és a mágneses térnek:
\begin{equation}
    \tilde{\bold{E}}(z, t) = \tilde{\bold{E}}_0 e^{i(kz - \omega t)} \qquad \tilde{\bold{B}}(z, t) = \tilde{\bold{B}}_0 e^{i(kz - \omega t)}
\end{equation}
Ahol $\tilde{E}_0$ és $\tilde{B}_0$ a komplex amplitúdók. Temészetesen a fizikai terek ezek valós részei lesznek. Az $\omega$ pedig a frekvencia ($\omega = ck$).
\begin{figure}[H]
    \centering
    \includegraphics[width=0.5\textwidth]{imgs/8-tetel/planaris_hullam.png}
    \caption{Planáris elektromágneses hullám vákuumban}
\end{figure}
Az egyenletek amiket kaptunk $\bold{E}$ és $\bold{B}$ terekre, a Maxwell-egyenletekből vezettük le, de ettől még az eredménynek teljesítenie kell a hullámegyenletet. De a Maxwell-egynlet is ad megközéseket,
mégpedig, mivel $\boldsymbol{\nabla} \cdot \bold{E} = 0$ és $\boldsymbol{\nabla} \cdot \bold{B} = 0$, ezért a tereknek divergenciamentesnek kell lenniük. Ez a feltétel a következőképpen írható fel a monokromatikus hullámokra :
\begin{equation}
    (\tilde{E}_0)_z = 0 \qquad (\tilde{B}_0)_z = 0
\end{equation}
Tehát az elektromos és mágneses tereknek merrőlegesnek kell lennie egymásnak, vagyis az elektromágneses hullámok \textbf{transzverzális hullámok}.
Továbbá a Faraday-törvény ($\boldsymbol{\nabla} \times \bold{E} = -\frac{\partial \bold{B}}{\partial t}$) feltételezi, hogy az elektromos és a mágneses terek amplitúdói között a következő kapcsolat áll fenn:
\begin{equation}
    -k(\tilde{E}_0)_y = i \omega (\tilde{B}_0)_x \qquad k(\tilde{E}_0)_x = i \omega (\tilde{B}_0)_y
\end{equation}
Vagy kompaktabb formában:
\begin{equation}
    \bold{B}_0 = \frac{k}{\omega} (\hat{\bold{z}} \times \tilde{\bold{E}}_0)
\end{equation}
Ezekszerint az elektromos és mágneses terek amplitúdói is merőlegesek egymásra, és a hullám terjedési irányára is. Tehát igazából csak egy nem nulla komponense lehet mindegyiknek. Emelett fázisban is Vannak
és a valós amplitúdójuk is kapcsolódik:
\begin{equation}
    B_0 = \frac{k}{\omega} E_0 = \frac{1}{c} E_0
\end{equation}
A Maxwell-Ampère törvény nem ad újabb megkötést, csak reprrodukálja az amplitúdók merőlegességét.

Tehát, ha például az elektromos tér ($\bold{E}$) az $x$ irányba mutat és a mágneses tér ($\bold{B}$) az $y$ irányba mutat, akkor a terek alakjai:
\begin{equation}
    \tilde{\bold{E}}(z, t) = \tilde{E}_0 e^{i(kz - \omega t )} \hat{\bold{x}} \qquad \tilde{\bold{B}}(z, t) = \frac{1}{c} \tilde{E}_0 e^{i(kz - \omega t)} \hat{\bold{y}}
\end{equation}
A fizikai terek pedig a valós részek:
\begin{equation}
    \boxed{\bold{E}(z, t) = E_0 \cos(kz - \omega t + \phi) \hat{\bold{x}} \qquad \bold{B}(z, t) = \frac{E_0}{c} \cos(kz - \omega t + \phi) \hat{\bold{y}}}
\end{equation}
Ez a monokromatikus planáris hullám lényege. Ezt a polarizáltságot konvenció alapján x irányú polarizációnak nevezzük.
\begin{figure}[H]
    \centering
    \includegraphics[width=0.5\textwidth]{imgs/8-tetel/x_polarizalt_hullam.png}
    \caption{x irányú polarizált elektromágneses hullám vákuumban}
\end{figure}
Természetesen nincsen semmi különleges a $z$ irányban, lehet általánosítani tetszőleges irányú terjedésre:
\begin{equation}
    \boxed{
    \begin{aligned}
        \bold{E}(\bold{r}, t) &= E_0 e^{i(\bold{k} \cdot \bold{r} - \omega t)} \hat{\bold{n}} \\
        \bold{B}(\bold{r}, t) &= \frac{E_0}{c} e^{i(\bold{k} \cdot \bold{r} - \omega t)} (\hat{\bold{k}} \times \hat{\bold{n}}) = \frac{1}{c} (\hat{\bold{k}} \times \bold{E} )
    \end{aligned}}
\end{equation}
Ahol $\hat{\bold{n}}$ a polarizációs vektor, ami merőleges a hullám terjedési irányára:
\begin{equation}
    \hat{\bold{n}} \cdot \hat{\bold{k}} = 0
\end{equation}
A mágneses tér pedig egyértelműen merőleges az elektromos térre a keresztszorzat miatt. A valós terek pedig a valós részek:
\begin{equation}
    \begin{aligned}
        \bold{E}(\bold{r}, t) &= E_0 \cos(\bold{k} \cdot \bold{r} - \omega t + \phi) \hat{\bold{n}} \\
        \bold{B}(\bold{r}, t) &= \frac{E_0}{c} \cos(\bold{k} \cdot \bold{r} - \omega t + \phi) (\hat{\bold{k}} \times \hat{\bold{n}})
    \end{aligned}
\end{equation}

\subsubsection{Elektromágneses hullámok energia és impulzusa}
Az elektromágneses tér energiája:
\begin{equation}
    u = \frac{1}{2} \left( \epsilon_0 E^2 + \frac{1}{\mu_0} B^2 \right)
\end{equation}
Ahol monokromatikus hullámok esetében:
\begin{equation}
    E = E_0 \cos(kz - \omega t + \phi) \qquad B = \frac{E_0}{c} \cos(kz - \omega t + \phi)
\end{equation}
Tehát:
\begin{equation}
    B^2 = \frac{E_0^2}{c^2} E^2 = \mu_0 \epsilon_0 E^2
\end{equation}
Vagyis az elektromos és a mágneses tér energiája megegyezik monokromatikus hullámok esetében:
\begin{equation}
    u = \epsilon_0 E^2 = \epsilon_0 E_0^2 \cos^2(kz - \omega t + \phi)
\end{equation}
Az enegia fluxus sűrűséget a Poynting-vektorral írhatjuk le:
\begin{equation}
    \bold{S} = \frac{1}{\mu_0} (\bold{E} \times \bold{B})
\end{equation}
Monokromatikus hullámok esetében:
\begin{equation}
    \bold{S} = c\epsilon_0 E^2 \cos^2(kz - \omega t + \phi) \hat{\bold{z}} = cu \hat{\bold{z}}
\end{equation}
A Poynting vektor tehát az energia sűrűség és a hullámsebesség szorzata, ami logikus, hiszen az energia fluxus sűrűség azt jelenti, hogy egységnyi területen egységnyi idő alatt mennyi energia halad át.
Tehát egy egységnyi területen egységnyi idő alatt $uc$ energiát szállít át a hullám.

De nem csak energiájuk van az elektromágneses hullámoknak, hanem impulzusuk is. Az elektromágneses hullámok impulzus sűrűsége a következőképpen adódik:
\begin{equation}
    \bold{g} = \frac{\bold{S}}{c^2}
\end{equation}
Monokromatikus hullámok esetében:
\begin{equation}
    \bold{g} = \frac{1}{c}\epsilon_0 E_0^2 \cos^2(kz - \omega t + \phi) \hat{\bold{z}} = \frac{u}{c} \hat{\bold{z}}
\end{equation}
De a fény esetében a hullámhossz túl kicsi és a frekvencia túl nagy, ezért minden makroszkopikus mérés sok hullámot fog magába foglalni. ezért nem a pillanatnyi értékekkel dolgozunk,
hanem az időátlagolt értékekkel. Az időátlagolt energia sűrűség és impulzus sűrűség a következőképpen adódik:
\begin{equation}
    \langle u \rangle = \frac{1}{2} \epsilon_0 E_0^2 \qquad \langle \bold{g} \rangle = \frac{1}{2c} \epsilon_0 E_0^2 \hat{\bold{z}} \qquad \langle \bold{S} \rangle = \frac{1}{2} c \epsilon_0 E_0^2 \hat{\bold{z}}
\end{equation}
Ahol $\langle ... \rangle$ az időátlagot jelöli. Az átlagos teljesítmény sűrűség egységnyi területre (intenzitás) pedig:
\begin{equation}
    I \equiv |\langle \bold{S} \rangle| = \frac{1}{2} c \epsilon_0 E_0^2
\end{equation}
Ha a fény egy tökéletes elnyelőbe ütközik, akkorr az impulzusát átadja, amit \textbf{sugárzási nyomásnak} nevezünk. Ez a nyomás az intenzitás és a fénysebesség hányadosa:
\begin{equation}
    P = \frac{1}{A} \frac{\Delta p}{\Delta t} = \frac{1}{2} \epsilon_0 E_0^2 =\frac{I}{c}
\end{equation}
Egy tökéletes visszaverő felületen pedig 2x ekkora mert az impulzus nem csak elnyelődik, hanem az ellentétes irányba vált át.

\subsubsection{Elektomágneses hullámok anyagban}
\subsubsection*{dielektrikum}
Anyagban, ha nincsenek szabad töltések és áramok, akkor a Maxwell-egyenletek a következőképpen alakulnak:
\begin{equation}
    \begin{aligned}
        \boldsymbol{\nabla} \cdot \bold{D} &= 0 \qquad \boldsymbol{\nabla} \times \bold{E} = -\frac{\partial \bold{B}}{\partial t} \\
        \boldsymbol{\nabla} \cdot \bold{B} &= 0 \qquad \boldsymbol{\nabla} \times \bold{H} = \frac{\partial \bold{D}}{\partial t} 
    \end{aligned}
\end{equation}
Ha a közeg lineáris (vagyis a polarizáció és a mágnesezhetőség arányos a térerősséggel), akkor:
\begin{equation}
    \bold{D} = \epsilon \bold{E} \qquad \bold{B} = \frac{1}{\mu} \bold{H}
\end{equation}
Ahol $\epsilon$ a közeg permittivitása és $\mu$ a közeg permeabilitása. Ha pedig homogén is (vagyis a permittivitás és permeabilitás nem függ a helytől), akkor a Maxwell-egyenletek tovább egyszerűsödnek:
\begin{equation}
    \begin{aligned}
        \boldsymbol{\nabla} \cdot \bold{E} &= 0 \qquad \boldsymbol{\nabla} \times \bold{E} = -\frac{\partial \bold{B}}{\partial t} \\
        \boldsymbol{\nabla} \cdot \bold{B} &= 0 \qquad \boldsymbol{\nabla} \times \bold{B} = \mu \epsilon \frac{\partial \bold{E}}{\partial t} 
    \end{aligned}
\end{equation}
Amik a vákuumbeli egyenletektől csak annyiban térnek el, hogy a vákuumbeli $\mu_0$ és $\epsilon_0$ helyett a közegbeli $\mu$ és $\epsilon$ szerepel. Innen ugyanúgy levezethető a hullámegyenlet:
\begin{equation}
    \boxed{\boldsymbol{\nabla}^2 \bold{E} = \mu \epsilon \frac{\partial^2 \bold{E}}{\partial t^2}} \qquad \boxed{\boldsymbol{\nabla}^2 \bold{B} = \mu \epsilon \frac{\partial^2 \bold{B}}{\partial t^2}}
\end{equation}
Ami azt jelenti, hogy az elektromágneses hullámoknak vannak hullámszerű megoldásai anyagban is. Az elektromágneses hullámok terjedési sebessége anyagban pedig a következőképpen adódik:
\begin{equation}
    v = \frac{1}{\sqrt{\mu \epsilon}} = \frac{c}{n} \qquad \text{ahol } n = \sqrt{\frac{\mu \epsilon}{\mu_0 \epsilon_0}}
\end{equation}
Ahol $n$ a közeg \textbf{törésmutatója}. Mivel a legtöbb anyagnak $\mu \approx \mu_0$, ezért a törésmutató főként a permittivitástól függ:
\begin{equation}
    n \approx \sqrt{\frac{\epsilon}{\epsilon_0}} \approx \sqrt(\epsilon_r)
\end{equation}
Ahol $\epsilon_r$ a közeg \textbf{relatív permittivitása}. Mivel a legtöbb anyagnak $\epsilon > \epsilon_0$, ezért a törésmutató általában nagyobb mint 1, így az elektromágneses hullámok anyagban lassabban terjednek mint vákuumban.

A többi megoldásunk is használható itt, az energia sűrűségnél például csak az $\epsilon_0$ helyett $\epsilon$-t kell használni és a mágneses térnél a $\mu_0$ helyett $\mu$-t:
\begin{equation}
    u = \frac{1}{2} \left( \epsilon E^2 + \frac{1}{\mu} B^2 \right)
\end{equation}
A Poynting-vektor is hasonlóan alakul:
\begin{equation}
    \bold{S} = \frac{1}{\mu} (\bold{E} \times \bold{B})
\end{equation}
Monokromatikus hullámoknál a frekvencia és a hullámszám függnek egymástól $\omega = vk = \frac{c}{n} k$, és a $\bold{B}$ amplitúdója $1/v$-vel lesz kisebb az $\bold{E}$ amplitúdójánál:
\begin{equation}
    B_0 = \frac{k}{\omega} E_0 = \frac{1}{v} E_0 = \frac{n}{c} E_0
\end{equation}
Az intenzitás pedig:
\begin{equation}
    I = \frac{1}{2} v \epsilon E_0^2 = \frac{1}{2} \frac{c}{n} \epsilon E_0^2
\end{equation}
A határfeltételek is nagyon hasonlóan fognak alakulni:
\begin{equation}
    \begin{aligned}
        \epsilon_1 E_1^\perp &= \epsilon_2 E_2^\perp \qquad E_{1}^\parallel = E_{2}^\parallel \\
        B_1^\perp &= B_2^\perp \qquad \frac{1}{\mu_1} B_{1}^\parallel = \frac{1}{\mu_2} B_{2}^\parallel
    \end{aligned}
\end{equation}

\subsubsection*{Vezetők}
A vezetők esetében már nem igaz, hogy nincsenek szabad töltések és áramok, hiszen a vezetőkben szabadon mozoghatnak a töltések. Ezért az áramsűrűség Ohm törvénye alapján a következőképpen írható fel:
\begin{equation}
    \bold{J}_f = \sigma \bold{E}
\end{equation}
Ahol $\sigma$ a vezető \textbf{vezetőképessége}. Ezt az áramsűrűséget be kell helyettesíteni a Maxwell-Ampère törvénybe:
\begin{equation}
    \boldsymbol{\nabla} \times \bold{B} = \mu \left( \bold{J}_f + \epsilon \frac{\partial \bold{E}}{\partial t} \right) = \mu \left( \sigma \bold{E} + \epsilon \frac{\partial \bold{E}}{\partial t} \right)
\end{equation}
Tehát a Maxwell-egyenletek a vezetőben a következőképpen alakulnak:
\begin{equation}
    \begin{aligned}
        \boldsymbol{\nabla} \cdot \bold{E} &= \frac{1}{\epsilon} \rho_f \qquad \boldsymbol{\nabla} \times \bold{E} = -\frac{\partial \bold{B}}{\partial t} \\
        \boldsymbol{\nabla} \cdot \bold{B} &= 0 \qquad \boldsymbol{\nabla} \times \bold{B} = \mu \left( \sigma \bold{E} + \epsilon \frac{\partial \bold{E}}{\partial t} \right)
    \end{aligned}
\end{equation}
A kontinuitási egyenlet a szabad töltésre:
\begin{equation}
    \boldsymbol{\nabla} \cdot \bold{J}_f = - \frac{\partial \rho_f}{\partial t}
\end{equation}
Ez pedig az Ohm-törvénnyel és a Gauss törvénnyel együtt a következőképpen alakul:
\begin{equation}
    \frac{\partial \rho_f}{\partial t} = -\sigma (\boldsymbol{\nabla} \cdot \bold{E}) = - \frac{\sigma}{\epsilon} \rho_f
\end{equation}
Ami egy elsőrendű differenciálegyenlet $\rho_f$-re, aminek a megoldása, amennyiben homogén közegben vagyunk:
\begin{equation}
    \rho_f(t) = \rho_f(0) e^{-\frac{\sigma}{\epsilon} t}
\end{equation}
Tehát a vezetőben a szabad töltés $\rho_f(0)$ exponenciális sebességgel szétoszlik, a karakterisztikus $\tau \equiv \frac{\epsilon}{\sigma}$ idővel.
Ez megerősíti azt a tudásunkat, hogy ha egy szabad töltést teszünk egy vezetőre, akkor az ki fog áramlani a szélekre. A $\tau$-t használhatjuk egyfajta mérőszámként arra, hogy
megmondjuk mennyire "jó" egy vezető. Egy ideális vezető esetében $\sigma \to \infty$, így $\tau \to 0$, vagyis a töltések azonnal eltűnnek a vezetőből. Egy "jó" vezető esetében pedig $\sigma$ nagy, és $\tau \ll 1/\omega$.
De ez a tranziens viselkedés nem érdekel most minket, tehát tegyük fel, hogy várunk elég időt, hogy a szabad töltések már eltűnjenek a vezetőkről. Ekkor a Maxwell-egyenletek a következőképpen alakulnak:
\begin{equation}
    \begin{aligned}
        \boldsymbol{\nabla} \cdot \bold{E} &= 0 \qquad \boldsymbol{\nabla} \times \bold{E} = -\frac{\partial \bold{B}}{\partial t} \\
        \boldsymbol{\nabla} \cdot \bold{B} &= 0 \qquad \boldsymbol{\nabla} \times \bold{B} = \mu \left( \sigma \bold{E} + \epsilon \frac{\partial \bold{E}}{\partial t} \right)
    \end{aligned}
\end{equation}
Ez eltér a nem vezető anyagokban lévő Maxwell-egyenletektől az extra $\mu \sigma \bold{E}$ tag miatt. Ezt a tagot figyelembe véve a hullámegyenlet a következőképpen alakul:
\begin{equation}
    \boldsymbol{\nabla}^2 \bold{E} = \mu \epsilon \frac{\partial^2 \bold{E}}{\partial t^2} + \mu \sigma \frac{\partial \bold{E}}{\partial t} \qquad \boldsymbol{\nabla}^2 \bold{B} = \mu \epsilon \frac{\partial^2 \bold{B}}{\partial t^2} + \mu \sigma \frac{\partial \bold{B}}{\partial t}
\end{equation}
Ezek az egyenletek továbbra is engednek planáris hullám megoldásokat:
\begin{equation}
    \tilde{\bold{E}}(z, t) = \tilde{\bold{E}}_0 e^{i(\tilde{k}z - \omega t)} \qquad \tilde{\bold{B}}(z, t) = \tilde{\bold{B}}_0 e^{i(\tilde{k}z - \omega t)}
\end{equation}
De a hullámszám most már komplex lesz a vezetőben:
\begin{equation}
    \tilde{k}^2 = \mu \epsilon \omega^2 + i \mu \sigma \omega
\end{equation}
Ahol $\tilde{k} = k + i \kappa$:
\begin{equation}
    k \equiv \omega \sqrt{\frac{\mu \epsilon}{2}} \left[ \sqrt{1 + \left( \frac{\sigma}{\epsilon \omega} \right)^2} + 1 \right]^{1/2} \qquad \kappa \equiv \omega \sqrt{\frac{\mu \epsilon}{2}} \left[ \sqrt{1 + \left( \frac{\sigma}{\epsilon \omega} \right)^2} - 1 \right]^{1/2}
\end{equation}
Ahol $k$ a hullámszám valós része, és $\kappa$ az imaginárius része. Az imaginárius része a hullámnak elnyelődését írja le a vezetőben:
\begin{equation}
    \tilde{\bold{E}}(z, t) = \tilde{E}_0 e^{-\kappa z} e^{i(kz - \omega t)} \qquad \tilde{\bold{B}}(z, t) = \tilde{B}_0 e^{-\kappa z} e^{i(kz - \omega t)}
\end{equation}
Az a távolság amit a hullám megtesz a vezetőben, mire az intenzitása $1/e$-ed részére csökken (kb. $1/3$), a \textbf{skin-mélység} (behatolási mélység):
\begin{equation}
    d \equiv \frac{1}{\kappa}
\end{equation}
A $\tilde{k}$ valós része pedig továbbra is meghatározza a frekvenciát, a terjedési sebességet és a refrakciós indexet:
\begin{equation}
    \lambda = \frac{2\pi}{k} \qquad v = \frac{\omega}{k} \qquad n = \frac{ck}{\omega}
\end{equation}

Az elnyelődő hullámok kielégítik a módosított hullámegyenletet bármilyen $\tilde{E}_0$ és $\tilde{B}_0$-ra, de a Maxwell-egyenletek további megkötéseket szabnak az amplitúdókra és a fázisokra.
A Gauss és a mágneses Gauss törvény miatt csak egy irányú komponense lehet mindkét térnek. Emiatt akár meg is választhatjuk a teret úgy, hogy az elektromos tér az $x$ irányba legyen polarizált és a hullám a $z$ irányba terjedjen:
\begin{equation}
    \tilde{\bold{E}}(z, t) = \tilde{E}_0 e^{-\kappa z} e^{i(kz - \omega t)} \hat{\bold{x}}
\end{equation}
A Faraday-törvényből pedig:
\begin{equation}
    \tilde{\bold{B}}(z, t) = \frac{\tilde{k}}{\omega} \tilde{E}_0 e^{-\kappa z} e^{i(kz - \omega t)} \hat{\bold{y}}
\end{equation}
A Maxwell-Ampère törvény pedig nem ad újabb megkötést. Ez megint kimondja, hogy az elektromos és mágneses terek merőlegesek egymásra és a hullám terjedési irányára is.

Mint minden komplex szám, $\tilde{k}$ is felírható poláris alakban:
\begin{equation}
    \tilde{k} = K e^{i \phi} \qquad \text{ahol } \phi = \tan^{-1} \left( \frac{\kappa}{k} \right)
\end{equation}
Ahol:
\begin{equation}
    K = |\tilde{k}| = \sqrt{k^2 + \kappa^2} = \omega \sqrt{\epsilon \mu \sqrt{1 + \left( \frac{\sigma}{\epsilon \omega} \right)^2}}
\end{equation}
Az amplitúdók pedig a $\tilde{\bold{E}}$ és $\tilde{\bold{B}}$ függvényei alapján kacsolódnak egymáshoz:
\begin{equation}
    \tilde{E}_0 = E_0 e^{i \phi_E} \qquad \tilde{B}_0 = B_0 e^{i \phi_B} \quad \rightarrow \quad B_0 e^{i \phi_B} = \frac{K e^{i \phi}}{\omega} E_0 e^{i\phi_E}
\end{equation}
\begin{figure}[H]
    \centering
    \includegraphics[width=0.5\textwidth]{imgs/8-tetel/hullam_vezeto.png}
    \caption{Elektromágneses hullám vezetőben}
\end{figure}
Látható, hogy a vezetőben az elektromos és mágneses terek között fáziskülönbség van:
\begin{equation}
    \Delta \phi = \phi_B - \phi_E = \phi = \tan^{-1} \left( \frac{\kappa}{k} \right)
\end{equation}
Ez azt jelenti, hogy a vezetőben a mágneses tér le van maradva az elektromos mögött. A valós aplitúdók $\bold{E}$ és $\bold{B}$ szintén függnek egymástól:
\begin{equation}
    \frac{B_0}{E_0} = \frac{K}{\omega} = \sqrt{\mu \epsilon \sqrt{1 + \left( \frac{\sigma}{\epsilon \omega} \right)^2}}
\end{equation}
Tehát a valós elektromos és mágneses terek pedig:
\begin{equation}
    \boxed{
    \begin{aligned}
        \bold{E}(z, t) &= E_0 e^{-\kappa z} \cos(kz - \omega t + \phi_E) \hat{\bold{x}} \\
        \bold{B}(z, t) &= B_0 e^{-\kappa z} \cos(kz - \omega t + \phi_E + \phi) \hat{\bold{y}}
    \end{aligned}}
\end{equation}

\subsubsection*{transzmisszió és reflexió}
Vegyünk egy határfelületet két különböző \textbf{dielektrikus} anyag között, amit egy felület reprezentál az $x-y$ síkban.
A sugár utazzon a $z$ irányba, $x$ irányba polarizálva, a bal oldalról érkezve ($z < 0$).
\begin{figure}[H]
    \centering
    \includegraphics[width=0.5\textwidth]{imgs/8-tetel/reflekcio_normal.png}
    \caption{Elektromágneses hullám határfelületen}
\end{figure}
Ekkor a beérkező hullámok egyenletei:
\begin{equation}
    \tilde{\bold{E}}_I(z, t) = \tilde{E}_{0I} e^{i(k_1 z - \omega t)} \hat{\bold{x}} \qquad \tilde{\bold{B}}_I(z, t) = \frac{\tilde{E}_{0I}}{v_1} e^{i(k_1 z - \omega t)} \hat{\bold{y}}
\end{equation}
Ekkor a visszavert hullámok egyenletei:
\begin{equation}
    \tilde{\bold{E}}_R(z, t) = \tilde{E}_{0R} e^{i(-k_1 z - \omega t)} \hat{\bold{x}} \qquad \tilde{\bold{B}}_R(z, t) = -\frac{\tilde{E}_{0R}}{v_1} e^{i(-k_1 z - \omega t)} \hat{\bold{y}}
\end{equation}
Ami visszafelé utazik az 1. közegben. Végül a továbbított hullámok egyenletei a 2. közegben:
\begin{equation}
    \tilde{\bold{E}}_T(z, t) = \tilde{E}_{0T} e^{i(k_2 z - \omega t)} \hat{\bold{x}} \qquad \tilde{\bold{B}}_T(z, t) = \frac{\tilde{E}_{0T}}{v_2} e^{i(k_2 z - \omega t)} \hat{\bold{y}}
\end{equation}
Ahol $k_1 = \frac{\omega}{v_1}$ és $k_2 = \frac{\omega}{v_2}$. Ahol $v_1$ és $v_2$ a hullámok sebessége az 1. és 2. közegben.Vegyük észre, hogy a $\tilde{\bold{B}}_r$-nél bejön egy negatív előjel,
mivel a $\bold{B}$ értéke függ a hullám terjedési irányától is. A $z = 0$ pontban a kombinált tereknek a bal oldalon kapcsolódniuk kell a jobb oldalon lévő terekhez a határfeltételek alapján.
Az első két határfeltétel triviálisan teljesül:
\begin{equation}
    \epsilon_1 E_{I}^\perp + \epsilon_1 E_{R}^\perp = \epsilon_2 E_{T}^\perp \quad \rightarrow \quad \epsilon_1 \tilde{E}_{0I} + \epsilon_1 \tilde{E}_{0R} = \epsilon_2 \tilde{E}_{0T}
\end{equation}
Viszont a harmadik megköveteli, hogy:
\begin{equation}
    \tilde{E}_{0I} + \tilde{E}_{0R} = \tilde{E}_{0T}
\end{equation}
A negyedik határfeltétel pedig:
\begin{equation}
    \frac{1}{\mu_1} \left( \frac{\tilde{E}_{0I}}{v_1} - \frac{\tilde{E}_{0R}}{v_1} \right) = \frac{1}{\mu_2} \frac{\tilde{E}_{0T}}{v_2}
\end{equation}
Vagy átrendezve:
\begin{equation}
    \tilde{E}_{0I} - \tilde{E}_{0R} = \frac{\mu_1 v_1}{\mu_2 v_2} \tilde{E}_{0T}
\end{equation}
Ahol bevezethetjük a $\beta$ paramétert:
\begin{equation}
    \beta \equiv \frac{\mu_1 v_1}{\mu_2 v_2} = \frac{\mu_1 n_2}{\mu_2 n_1}
\end{equation}
Most már négy egyenletünk van négy ismeretlennel ($\tilde{E}_{0I}, \tilde{E}_{0R}, \tilde{E}_{0T}$). Ezek megoldásával megkaphatjuk a visszavert és továbbított hullámok amplitúdóit a beérkező hullám amplitúdójához képest:
\begin{equation}
    \tilde{E}_{0R} = \left(\frac{1 - \beta}{1 + \beta}\right) \tilde{E}_{0I} \qquad \tilde{E}_{0T} = \left(\frac{2}{1 + \beta}\right) \tilde{E}_{0I}
\end{equation}
Ezek a megoldások nagyon hasonlóak azokhoz amiket a rezgetett köteleken kaptunk. Ha a $\mu$ értéke közel van a vákuumbeli értékéhez (ami a legtöbb közegben igaz), akkor $\beta = v_1/v_2$, vagyis:
\begin{equation}
    \tilde{E}_{0R} = \left(\frac{v_2 - v_1}{v_2 + v_1}\right) \tilde{E}_{0I} \qquad \tilde{E}_{0T} = \left(\frac{2 v_2}{v_2 + v_1}\right) \tilde{E}_{0I}
\end{equation}
Ami teljesen megegyezik a rezgetett kötelek eredményével a sebességek helyettesítésével. Itt tehát ugyan úgy igaz, hogy a fázisok azonosak amíg $v_2 > v_1$, de ha $v_2 < v_1$, akkor a visszavert hullám fázisa eltolódik $\pi$-vel.
A valós amplitúdók pedig:
\begin{equation}
    E_{0R} = \left|\frac{v_2 - v_1}{v_2 + v_1}\right| E_{0I} \qquad E_{0T} = \left(\frac{2 v_2}{v_2 + v_1}\right) E_{0I}
\end{equation}
Mivel a refrakciós index $n = c/v$, ezért felírhatjuk a visszaverési és transzmissziós együtthatókat is a refrakciós indexek segítségével is:
\begin{equation}
    E_{0R} = \left|\frac{n_1 - n_2}{n_1 + n_2}\right| E_{0I} \qquad E_{0T} = \left(\frac{2 n_1}{n_1 + n_2}\right) E_{0I}
\end{equation}
De kérdés, hogy mennyi energia tükröződik vissza és mennyi megy át a határfelületen. Az intenzitás megadható az elektromos tér amplitúdójával:
\begin{equation}
    I = \frac{1}{2} v \epsilon E_0^2
\end{equation}
És ha $\mu_1 = \mu_2 \approx \mu_0$ továbbra is fennáll, akkor a visszaverődő intenzitás és a teljes intenzitás aránya:
\begin{equation}
    R \equiv \frac{I_R}{I_I} = \left( \frac{E_{0R}}{E_{0I}} \right)^2 = \left(\frac{n_1 - n_2}{n_1 + n_2}\right)^2
\end{equation}
Ez a \textbf{reflexiós együttható}. A transzmissziós együttható pedig:
\begin{equation}
    T \equiv \frac{I_T}{I_I} = \frac{v_2 \epsilon_2 E_{0T}^2}{v_1 \epsilon_1 E_{0I}^2} = \frac{4 n_1 n_2}{(n_1 + n_2)^2}
\end{equation}
Ahol használtuk, hogy $n = c/v$ és $\epsilon \approx n^2 \epsilon_0$. Könnyen ellenőrizhető, hogy $R + T = 1$, vagyis az energia megmarad.

De mi van akkor, ha a beesési szög nem pont $90^\circ$? Ekkor a bejövő hullámegyenletek a következőképpen alakulnak:
\begin{equation}
    \tilde{\bold{E}}_I(\bold{r}, t) = \tilde{E}_{0I} e^{i(\bold{k}_1 \cdot \bold{r} - \omega t)} \hat{\bold{n}}_I \qquad \tilde{\bold{B}}_I(\bold{r}, t) = \frac{1}{v_1} (\hat{\bold{k}}_I \times \tilde{\bold{E}}_I)
\end{equation}
A visszaverődő hullámok pedig:
\begin{equation}
    \tilde{\bold{E}}_R(\bold{r}, t) = \tilde{E}_{0R} e^{i(\bold{k}_R \cdot \bold{r} - \omega t)} \hat{\bold{n}}_R \qquad \tilde{\bold{B}}_R(\bold{r}, t) = \frac{1}{v_1} (\hat{\bold{k}}_R \times \tilde{\bold{E}}_R)
\end{equation}
\begin{figure}[H]
    \centering
    \includegraphics[width=0.5\textwidth]{imgs/8-tetel/refleckio_szog.png}
    \caption{Elektromágneses hullám határfelületen nem normális beesés esetén}
\end{figure}
A továbbított hullámok pedig:
\begin{equation}
    \tilde{\bold{E}}_T(\bold{r}, t) = \tilde{E}_{0T} e^{i(\bold{k}_2 \cdot \bold{r} - \omega t)} \hat{\bold{n}}_T \qquad \tilde{\bold{B}}_T(\bold{r}, t) = \frac{1}{v_2} (\hat{\bold{k}}_T \times \tilde{\bold{E}}_T)
\end{equation}
Mindegyik hullámnak azonos a $\omega$ frekvenciája, amit csak a forrás határoz meg. A hullámszámok viszont különbözőek lehetnek, de függnek egymástól:
\begin{equation}
    k_1 v = k_R v_1 = k_T v_2 = \omega \quad \text{vagy} \quad k_1 = k_R = \frac{v_2}{v_1} k_T = \frac{n_1}{n_2} k_T
\end{equation}
A terek kapcsolódását megint a határfeltételek adják meg és ezeknek az egyenleteknek van egy általános struktúrája:
\begin{equation}
    ()e^{i(\bold{k}_1 \cdot \bold{r} - \omega t)} + ()e^{i(\bold{k}_R \cdot \bold{r} - \omega t)} = ()e^{i(\bold{k}_2 \cdot \bold{r} - \omega t)} \qquad z = 0\text{-nál}
\end{equation}
Ebből az alakból az a fontos tanulság, hogy az $x, y, t$ függés mind az exponensbe van "koncentrálva". Mivel pedig a határfeltételeknek a felület minden pontján igaznak kell lennie,
és minden időpontban teljesülniük kell, ezért az exponenseknek egyenlőnek kell lennie $z=0$-ban! Különben egy kis változás, mondjuk $x$-ben egyből tönkretenné az egyenlőséget.
Természetesen az idő komponensek már eleve egyenlőek, és a tér komponensekre igaz, hogy:
\begin{equation}
    \bold{k}_1 \cdot \bold{r} = \bold{k}_R \cdot \bold{r} = \bold{k}_2 \cdot \bold{r} \qquad \text{ha } z = 0
\end{equation}
Vagy explicitebben:
\begin{equation}
    x(k_I)_x + y(k_I)_y = x(k_R)_x + y(k_R)_y = x(k_T)_x + y(k_T)_y \qquad \text{ha } z = 0
\end{equation}
Minden $x$ és $y$ értékre. Ez csak akkor teljesülhet, ha a hullámszámok $x$ és $y$ komponensei megegyeznek:
\begin{equation}
    \begin{cases}
        (k_I)_x = (k_R)_x = (k_T)_x \quad \text{ha, } x = 0\\
        (k_I)_y = (k_R)_y = (k_T)_y \quad \text{ha, } y = 0
    \end{cases}
\end{equation}
Ezekből az egyenletekből pedig levezethetjük az alapvető törvényeit a hullámok visszaverődésére és törésére:
\begin{quote}
     \textbf{Első törvény}: A bejövő, visszavert és továbbított hullámok vektorai egy felületet alkotnak (amit \textbf{beesési síknak} nevezünk), és ami tartalmazza a felület normál vektorát is. 
     Miközben a hullámszám egyenlete $y = 0$-ban implikálja, hogy:
        \begin{equation}
            k_I sin \theta_I = k_R sin \theta_R = k_T sin \theta_T
        \end{equation}
\end{quote}
Ahol $\theta_I, \theta_R, \theta_T$ a bejövő, visszavert és továbbított hullámok beesési szögei a felület normáljához képest.
\begin{quote}
    \textbf{Második törvény}: A visszavert hullám beesési szöge megegyezik a bejövő hullám beesési szögével:
        \begin{equation}
            \theta_R = \theta_I
        \end{equation}
\end{quote}
Ezt nevezik \textbf{visszaverődési törvénynek}.
\begin{quote}
    \textbf{Harmadik törvény}: A továbbított hullám beesési szöge és a bejövő hullám beesési szöge között a következő kapcsolat áll fenn:
        \begin{equation}
            \frac{sin \theta_T}{sin \theta_I} = \frac{n_1}{n_2}
        \end{equation}
\end{quote}
Ezt nevezik \textbf{Snellius-Descartes-törvénynek}. Ez a három törvény a geometriai optika fundamentális egyenletei. Érdekesség, hogy nagyon kevés "elektrodinamika" került ezekbe az egyenletekbe,
semmilyen specifikus határfeltételt nem hoztunk még be, tehát ezek a törvények általánosabban is igazak a hullámokra, nem csak az elektromágnesesekre.

Most, hogy az exponenciális faktorokat rendbe tettük, hisz láttuk $z = 0$-ban kiesnek, a határfeltételek a következőképpen alakulnak:
\begin{equation}
    \begin{aligned}
        (i) \quad & \epsilon_1 (\tilde{E}_{0I} + \tilde{E}_{0R})_z = \epsilon_2 (\tilde{E}_{0T})_z \\
        (ii) \quad & (\tilde{\bold{B}}_{0I} + \tilde{\bold{B}}_{0R})_z = (\tilde{\bold{B}}_{0T})_z \\
        (iii) \quad & (\tilde{\bold{E}}_{0I} + \tilde{\bold{E}}_{0R})_z = (\tilde{\bold{E}}_{0T})_{x,y} \\
        (iv) \quad & \frac{1}{\mu_1} (\tilde{\bold{B}}_{0I} + \tilde{\bold{B}}_{0R})_{x,y} = \frac{1}{\mu_2} (\tilde{\bold{B}}_{0T})_{x,y}
    \end{aligned}
\end{equation}
Ahol $\tilde{\bold{B}}_0 = \frac{1}{v} (\hat{\bold{k}} \times \tilde{\bold{E}}_0)$. Az $(iii)$ és $(iv)$ egyenletek két egyenletet reprezentálnak $x$ és $y$ komponensekre külön-külön.
\begin{figure}[H]
    \centering
    \includegraphics[width=0.5\textwidth]{imgs/8-tetel/refrakcio.png}
    \caption{Beesési szög és törési szög}
\end{figure}

Tegyük fel, hogy a beeső hullám polarizációja a beesési síkban van ($x-y$ sík). Ekkor a továbbmenő és visszavert hullámok is ebben a síkban lesznek polarizálva.
Ekkor a $(i)$ határfeltétel:
\begin{equation}
    \epsilon_1 (- \tilde{E}_{0I} \sin \theta_I + \tilde{E}_{0R} \sin \theta_R) = \epsilon_2 (- \tilde{E}_{0T} \sin \theta_T)
\end{equation}
A $(ii)$ feltétel nem ad semmi érdekeset, és mivel a mágneses tereknek nincsen $z$ komponense, ezért $(iii)$:
\begin{equation}
    \tilde{E}_{0I} \cos \theta_I + \tilde{E}_{0R} \cos \theta_R = \tilde{E}_{0T} \cos \theta_T
\end{equation}
Végül a $(iv)$ határfeltétel:
\begin{equation}
    \frac{1}{\mu_1 v_1} \left( \tilde{E}_{0I} - \tilde{E}_{0R} \right) = \frac{1}{\mu_2 v_2} \tilde{E}_{0T}
\end{equation}
Ha felhasználjuk a visszaverődési és törési törvényeket, akkor ezek az egyenletek a következőképpen egyszerűsödnek:
\begin{equation}
    \tilde{E}_{0I} - \tilde{E}_{0R} = \beta \tilde{E}_{0T}
\end{equation}
Ahol (ahogy korábban is):
\begin{equation}
    \beta \equiv \frac{\mu_1 v_1}{\mu_2 v_2} = \frac{\mu_1 n_2}{\mu_2 n_1}
\end{equation}
És a $(iii)$ határfeltétel redukált alakjából pedig:
\begin{equation}
    \tilde{E}_{0I} + \tilde{E}_{0R} = \alpha \tilde{E}_{0T}
\end{equation}
Ahol:
\begin{equation}
    \alpha \equiv \frac{\cos \theta_T}{\cos \theta_I}
\end{equation}
Most már megoldhatjuk ezt a két egyenletet a visszavert és továbbított hullámok amplitúdóira a bejövő hullám amplitúdójához képest:
\begin{equation}
    \boxed{\tilde{E}_{0R} = \left( \frac{\alpha - \beta}{\alpha + \beta} \right) \tilde{E}_{0I} \qquad \tilde{E}_{0T} = \left( \frac{2}{\alpha + \beta} \right) \tilde{E}_{0I}}
\end{equation}
Ezek a \textbf{Fresnel-egyenletek} a beesési síkban polarizált hullámokra. Hasonlóan, ha a hullám polarizációja merőleges a beesési síkra (azaz $y$ irányba), akkor a Fresnel-egyenletek a következőképpen alakulnak:
\begin{equation}
    \boxed{\tilde{E}_{0R} = \left( \frac{1 - \beta \alpha}{1 + \beta \alpha} \right) \tilde{E}_{0I} \qquad \tilde{E}_{0T} = \left( \frac{2}{1 + \beta \alpha} \right) \tilde{E}_{0I}}
\end{equation}
Észrevehetjük, hogy az átmenő hullám mindig fázisban vannak a bejövő hullámmal, de a visszavert hullám pedig fázisban van, ha $\alpha > \beta$, de $\pi$-vel el van tolódva, ha $\alpha < \beta$.

A visszaverődési és transzmissziós együtthatók pedig a következőképpen alakulnak a beesési síkban polarizált hullámokra:
\begin{equation}
    R_\parallel \equiv \frac{I_R}{I_I} = \left( \frac{\alpha - \beta}{\alpha + \beta} \right)^2 \qquad T_\parallel \equiv \frac{I_T}{I_I} = \frac{4 \alpha \beta}{(\alpha + \beta)^2}
\end{equation}
Míg a merőleges polarizáció esetén:
\begin{equation}
    R_\perp \equiv \frac{I_R}{I_I} = \left( \frac{1 - \beta \alpha}{1 + \beta \alpha} \right)^2 \qquad T_\perp \equiv \frac{I_T}{I_I} = \frac{4 \beta \alpha}{(1 + \beta \alpha)^2}
\end{equation}

\subsubsection*{Visszaverődés Vezetőkön}
A határfeltételek megváltoznak egy kicsit, ha a szabad töltés és a szabad áram nem nulla a határfelületen. Ez akkor fordul elő, ha az egyik közeg egy jó vezető. Ebben az esetben a határfeltételek a következőképpen alakulnak:
\begin{equation}
    \begin{aligned}
        (i) \quad & \epsilon_1 E_1^{\perp} - \epsilon_2 E_2^{\perp} = \sigma_f \\
        (ii) \quad & B_1^\perp - B_2^\perp = 0 \\
        (iii) \quad & \bold{E}_1^\parallel - \bold{E}_2^\parallel = \bold{0} \\
        (iv) \quad & \frac{1}{\mu_1} \bold{B}_1^\parallel - \frac{1}{\mu_2} \bold{B}_2^\parallel = \mu_0 \bold{K}_f \times \hat{\bold{n}}
    \end{aligned}
\end{equation} 
Ahol $\sigma_f$ a felületi szabad töltéssűrűsé (nem összekeverendő a vezetőképességgel), $\bold{K}_f$ a felületi szabad áram sűrűség és $\hat{\bold{n}}$ a határfelület normálvektora (nem összekeverendő a polarizáció vektorral). 

Most tegyük fel, hogy egy felület határt képez két közeg között $x-y$ síkban, ahol az 1. közeg egy dielektrikum, a 2. közeg pedig egy jó vezető. Ekkor egy monokromatikus hullám x irányba polarizálva utazik
z irányban balról. Ekkor a bejövő hullámok egyenletei:
\begin{equation}
    \tilde{\bold{E}}_I(z, t) = \tilde{E}_{0I} e^{i(k_1 z - \omega t)} \hat{\bold{x}} \qquad \tilde{\bold{B}}_I(z, t) = \frac{\tilde{E}_{0I}}{v_1} e^{i(k_1 z - \omega t)} \hat{\bold{y}}
\end{equation}
A visszavert hullámok egyenletei:
\begin{equation}
    \tilde{\bold{E}}_R(z, t) = \tilde{E}_{0R} e^{i(-k_1 z - \omega t)} \hat{\bold{x}} \qquad \tilde{\bold{B}}_R(z, t) = -\frac{\tilde{E}_{0R}}{v_1} e^{i(-k_1 z - \omega t)} \hat{\bold{y}}
\end{equation}
A továbbmenő hullámok pedig a vezetőben:
\begin{equation}
    \tilde{\bold{E}}_T(z, t) = \tilde{E}_{0T} e^{i(\tilde{k}z - \omega t)} \hat{\bold{x}} \qquad \tilde{\bold{B}}_T(z, t) = \frac{\tilde{k}}{\omega} \tilde{E}_{0T} e^{i(\tilde{k}z - \omega t)} \hat{\bold{y}}
\end{equation}
Ahol $\tilde{k}$ a vezetőben lévő komplex hullámszám.

Most alkalmazzuk a határfeltételeket a $z = 0$ pontban. Mivel $E^{\perp} = 0$ és $B^\perp = 0$ ezért $(i)$ és $(ii)$ feltételek automatikusan teljesünek. A $(iii)$ határfeltétel:
\begin{equation}
    \tilde{E}_{0I} + \tilde{E}_{0R} = \tilde{E}_{0T}
\end{equation}
Míg a $(iv)$ határfeltétel:
\begin{equation}
    \frac{1}{\mu_1} \left( \frac{\tilde{E}_{0I}}{v_1} - \frac{\tilde{E}_{0R}}{v_1} \right) = \frac{1}{\mu_2} \frac{\tilde{k}}{\omega} \tilde{E}_{0T} = 0
\end{equation}
Vagy átrendezve:
\begin{equation}
    \tilde{E}_{0I} - \tilde{E}_{0R} = \tilde{\beta} \tilde{E}_{0T}
\end{equation}
Ahol:
\begin{equation}
    \tilde{\beta} \equiv \frac{\mu_1 v_1}{\mu_2 \omega} \tilde{k}_2
\end{equation}
Most már megoldhatjuk ezt a két egyenletet a visszavert és továbbított hullámok amplitúdóira a bejövő hullám amplitúdójához képest:
\begin{equation}
    \boxed{\tilde{E}_{0R} = \left( \frac{1 - \tilde{\beta}}{1 + \tilde{\beta}} \right) \tilde{E}_{0I} \qquad \tilde{E}_{0T} = \left( \frac{2}{1 + \tilde{\beta}} \right) \tilde{E}_{0I}}
\end{equation}
Ami rendkívül hasonló a dielektrikum-dielektrikum határfelületen kapott eredményhez, csak itt $\tilde{\beta}$ komplex szám.

Egy tökéletes vezetőben $\sigma = 0$ ezért $k_2 = \infty$ és így $\tilde{\beta} \rightarrow \infty$. Ekkor a visszaverődési együttható:
\begin{equation}
    \tilde{E}_{0R} = - \tilde{E}_{0I} \quad \text{és} \quad \tilde{E}_{0T} = 0 \qquad \text{vagyis} \qquad E_{0R} = E_{0I} \quad \text{és} \quad \phi_R = \phi_I + \pi
\end{equation}
Tehát egy tökéletes vezetőnél a hullám teljesen visszaverődik, és a visszavert hullám fázisa eltolódik $\pi$-vel a bejövő hullámhoz képest. Tehát a jó vezetők jó tükrök is.
Az otthon használt tükrökben se a dielektrikus üveg a tükröződő felület, az csak a struktúra megtartása miatt van ott, ami a lényegi "munkát" végzi az egy vékony réteg ezüst ami fel van festve rá.
Mivel az ezüst behatolási mélysége az optikai tartományba kevesebb mint $100 \: \text{\AA}$ ($1 \: \text{\AA} = 10^{-10} \: \text{m}$), ezért elég egy nagyon vékony réteg is a jó tükröződéshez.

\subsection{Diszperzió, csoport- és fázissebesség, Doppler-effektus}
Láttuk tehát, hogy az anyagban terjedő elektromágneses hullámok függnek az anyag három tulajdonságától:
\begin{itemize}
    \item dielektromos állandó $\epsilon$
    \item permeabilitás $\mu$
    \item vezetőképesség $\sigma$
\end{itemize}
Ezek az anyagtulajdonságok pedig valamilyen mértékben függnek a frekvenciától is. Például optikában jól ismert, hogy a törési index $n = \sqrt{\frac{\epsilon \mu}{\epsilon_0 \mu_0}} \approx \sqrt{\epsilon_r}$ a hullámhossz függvénye.
Ez azt jelenti, hogy különböző hullámhosszú (és ezáltal más színű) fény máshogy törik egy határfelületen. Ezt a jelenséget figyelhetjük meg a prizmánál is, ahol a különböző színű fények kombinációjából álló
beérkező fehér fény komponensei eltérő szögekben törnek meg, és emiatt szeparálódnak a prizmával való találkozás után. Ezt a jelenséget nevezzük \textbf{diszperziónak}. Ha pedig egy hullám sebessége függ a frekvenciájától, akkor közeget amiben ez a 
szituáció fennáll \textbf{diszperzív közegnek} nevezzük.

A különböző fekvenciájú hullámok nem csak különböző irányokba törnek meg, hanem különböző sebességgel is haladnak a közegben, ezért a hullám alakja, amiben több frekvencia is jelen van, megváltozik.
Egy erősen csúcsos hullám általában szétfolyik, és bár mindegyik szinuszos komponens a már ismert hullámsebességgel (vagy \textbf{fázissebességgel}) halad:
\begin{equation}
    v_f = \frac{\omega}{k} = \frac{c}{n(\omega)}
\end{equation}
A hullámok együttes alakja (hullámcsomag) sebessége más lesz. Ezt \textbf{csoportsebességnek} nevezzük, és a következőképpen definiáljuk:
\begin{equation}
    v_g = \frac{d\omega}{dk}
\end{equation}
Ez azt jelenti, hogy a hullámok egy "levelet" alkotnak ami a csoportsebességgel halad, miközben a benne lévő szinuszos komponensek a fázissebességgel haladnak. 

\begin{figure}[H]
    \centering
    \includegraphics[width=0.5\textwidth]{imgs/8-tetel/csoport_fazis_sebesseg.png}
    \caption{Csoportsebesség és fázissebesség}
\end{figure}
Ha például veszünk egy hullámot ami több különböző frekvenciából áll össze akkor tegyük fel, hogy a hullám a következő alakú:
\begin{equation}
    E(z, t) = \int_{-\infty}^{\infty} \tilde{E}(k) e^{i(kz - \omega(k)t)} dk
\end{equation}
Ahol a $\tilde{E}(k)$ a hullámcsomag spektruma, ami megadja, hogy az egyes hullámszámok milyen súllyal vannak jelen a hullámban. Tegyük fel, hogy a spektrum egy keskeny tartományban van középen
$k_0$ körül, tehát csak egy kis tartományban vannak jelen a hullámszámok. Ekkor a $\omega(k)$ függvényt ki tudjuk bővíteni Taylor-sorba $k_0$ körül:
\begin{equation}
    \omega(k) \approx \omega(k_0) + \left. \frac{d\omega}{dk} \right|_{k_0} (k - k_0) + \ldots
\end{equation}
Ha csak az első két tagot tartjuk meg, akkor a hullám a következőképpen alakul:
\begin{equation}
    E(z, t) \approx e^{i(k_0 z - \omega(k_0) t)} \int_{-\infty}^{\infty} \tilde{E}(k) e^{i(k - k_0)(z - v_g t)} dk
\end{equation}
Ahol bevezettük a csoportsebességet:
\begin{equation}
    v_g = \left. \frac{d\omega}{dk} \right|_{k_0}
\end{equation}
Ebből az egyenletből látható, hogy a hullám két részből áll: egy gyorsan oszcilláló rész, ami a fázissebességgel halad, és egy burkoló rész, ami a csoportsebességgel halad.

Érdekesség, de az energia, amit a hullámcsomag szállít nem a fázissebességgel, hanem a csoportsebességgel halad. Tehát előfordulhat, hogy a fázissebesség nagyobb, mint a fénysebesség.
Ez elsőre a relativitáselmélet megsértésének tűnhet, de valójában nem az, mivel a relativitáselmélet csak az információ terjedésére ad megkötést. Olyan dolgokról amik nem hordoznak magukban információt, a relativitáselmélet nem mond semmit.
És a fázissebesség önmagában nem tartalmaz információt. Egy hullám egy adott pontjának fázisa nem hordoz információt, csak a hullámcsomag alakja és annak változása hordoz információt.
Ezért nem sértjük meg a relativitáselméletet, ha a fázissebesség nagyobb, mint a fénysebesség.

\subsubsection{Doppler-Effektus}
Ha egy hullámforrás és egy megfigyelő egymáshoz képest mozognak, akkor a megfigyelő által mért frekvencia eltér a forrás által kibocsátott frekvenciától. Ezt a jelenséget nevezzük \textbf{Doppler-effektusnak}.
Ha a forrás és a megfigyelő közelednek egymáshoz, akkor a megfigyelt frekvencia nagyobb lesz, míg ha távolodnak egymástól, akkor kisebb lesz a megfigyelt frekvencia.
Matematikailag ez a következőképpen írható fel: 
\begin{equation}
    \nu' = \nu \left( \frac{v \pm v_{m}}{v \pm v_f} \right)
\end{equation}
Ahol $\nu$ a forrás által kibocsátott frekvencia, $v$ a hullám terjedési sebessége a közegben, $v_m$ a megfigyelő sebessége a közeghez képest (pozitív, ha a forrás felé mozog), és $v_f$ a forrás sebessége a közeghez képest (pozitív, ha a megfigyelő felé mozog).
Ha a forrás és a megfigyelő közelednek egymáshoz, akkor a képletben mindkét sebesség előjele pozitív lesz, és így a megfigyelt frekvencia nagyobb lesz, mint a kibocsátott frekvencia.
Ha pedig távolodnak egymástól, akkor mindkét sebesség előjele negatív lesz, és így a megfigyelt frekvencia kisebb lesz, mint a kibocsátott frekvencia.
Egy egyszerűbb esetet is megvizsgálhatunk, amikor a megfigyelő áll, és csak a forrás mozog egyenletes sebességgel $v$ a megfigyelő felé. Ekkor a megfigyelt frekvencia $\nu'$ a következőképpen alakul:
\begin{equation}
    \nu' = \nu \left( \frac{v}{v \pm v_f} \right)
\end{equation} 

Ahol $\nu$ a forrás által kibocsátott frekvencia. Ha a forrás és a megfigyelő távolodnak egymástól, akkor a képletben a $v$ előjele negatív lesz.
Ez a jelenség mind a mechanikai hullámokra, mind az elektromágneses hullámokra igaz. Az elektromágneses hullámok esetén azonban figyelembe kell venni a relativisztikus hatásokat is, mivel az elektromágneses hullámok sebessége mindig $c$ a vákuumban. Ha egy forrás egyenletes sebességgel $v$ mozog a megfigyelő felé, akkor a megfigyelt frekvencia $\nu'$ a következőképpen alakul:
\begin{equation}
    \nu' = \nu \sqrt{\frac{1 + \frac{v}{c}}{1 - \frac{v}{c}}}
\end{equation}
Ahol $\nu$ a forrás által kibocsátott frekvencia. Ha a forrás és a megfigyelő távolodnak egymástól, akkor a képletben a $v$ előjele negatív lesz.

A Doppler-effektusnak egyik legfontosabb alkalmazása a csillagászatban van, ahol a távoli csillagok és galaxisok mozgását tanulmányozzák a kibocsátott fény hullámhosszának változásán keresztül.
Ezt a jelenséget nevezik vöröseltolódásnak (ha a forrás távolodik) vagy kékeltolódásnak (ha a forrás közeledik). A vöröseltolódás mértéke segít meghatározni a csillagok és galaxisok sebességét és távolságát a Földtől.

\subsection{interferencia, diffrakció}
\subsubsection{Interferencia}
Az interferencia jelensége akkor fordul elő, amikor két vagy több hullám találkozik ugyanabban a térben, és ezek a hullámok kölcsönhatásba lépnek egymással.
Az interferencia eredményeként a hullámok amplitúdója megváltozik, és bizonyos helyeken erősödik, míg más helyeken gyengül vagy kioltódik. Ez a jelenség jól ismert például a vízfelszínen, ahol két hullámforrás által keltett hullámok találkoznak és interferálnak egymással.
Az interferencia jelenségét matematikailag a hullámok szuperpozíciójával írhatjuk le. Ha két hullám találkozik, akkor az eredő hullám amplitúdója az egyes hullámok amplitúdóinak összege lesz:
\begin{equation}
    E_{eredő} = E_1 + E_2
\end{equation}
Ahol $E_1$ és $E_2$ az egyes hullámok amplitúdói. Az interferencia eredményeként az eredő hullám amplitúdója függ a hullámok fáziskülönbségétől is. Ha a hullámok fázisban vannak, akkor az amplitúdók összeadódnak, és erősödik a hullám. Ha pedig a hullámok ellenfázisban vannak, akkor az amplitúdók kioltódnak, és gyengül a hullám.
Az interferencia jelensége számos területen fontos szerepet játszik, például az optikában, ahol a fényhullámok interferenciája révén jönnek létre különböző optikai jelenségek, mint például a Young-féle kettős rés kísérlet.
Az interferencia jelensége továbbá alapvető fontosságú a modern technológiákban is, például a lézerekben és a holográfiában.

\subsubsection*{Két-nyílás kísérlet}
A diffrakció szemléltetésének egyik legegyszerűbb módja a két-nyílás kísérlet ahol veszünk egy vékony falat amin két aprró nyílás van (a nyílás most legyen összemérhető a hullámhosszal).
Ekkor ha egy monokromatikus fényt bocsájtunk a nyílások felé, akkor a nyílások hullámforrásként fognak viselkedni, de mivel a nyílások $d$ távolságra vannak egymástól, ezért a két hullám el 
lesz tolódva egymáshoz képest. Ez azt okozza, hogy a hullámok egyes helyeken kioltják egymást, egyes helyeken pedig erősítik. Ez a jelenség jól látható a fal mögötti képernyőn, ahol egy interferencia mintázat jön létre.
\begin{figure}[H]
    \centering
    \includegraphics[width=0.5\textwidth]{imgs/8-tetel/kettos_nyilas.png}
    \caption{Két-nyílás kísérlet}
\end{figure}

Ekkor, a fény amplitúdója $u_p$ egy távoli $P$ pontban, $D$ távolságban:
\begin{equation}
    u_p = \frac{u_0}{r_1} e^{i(\omega t - k r_1)} + \frac{u_0}{r_2} e^{i(\omega t - k r_2)}
\end{equation}
Vegyük $r_1 - r_2 = d \sin \theta$, ahol $\theta$ a $P$ pont és a középvonal által bezárt szög. Ekkor ha feltesszük, hogy nagy távolságban $r_1 \approx r_2 = r$, akkor az intenzitás:
\begin{equation}
    I_p = 4 \left(\frac{u_0}{r}\right)^2 \cos^2 \left( \frac{kd \sin \theta}{2} \right)
\end{equation}

\subsubsection{diffrakció}
A diffrakció jelensége akkor fordul elő, amikor egy hullám találkozik egy akadállyal vagy réssel, és a hullám terjedési iránya megváltozik. A diffrakció eredményeként a hullám elhajlik az akadály körül, és új hullámfrontok jönnek létre. Ez a jelenség jól ismert például a vízfelszínen, ahol egy akadály elhelyezésekor a hullámok elhajlanak az akadály körül, és új hullámfrontok jönnek létre.
A diffrakció jelenségét matematikailag a Huygens-Fresnel elvvel írhatjuk le:
\begin{quote}
    \textbf{Huygens-Fresnel elv}: Minden pont egy hullámfronton úgy viselkedik, mint egy újabb hullámforrás, amelyből szférikus hullámok indulnak ki. Az új hullámfrontot ezeknek a szférikus hullámoknak az interferenciája hozza létre.
\end{quote}
A diffrakció jelensége számos területen fontos szerepet játszik, például az optikában, ahol a fényhullámok diffrakciója révén jönnek létre különböző optikai jelenségek, mint például a rácsokon való elhajlás.
A diffrakció jelensége továbbá alapvető fontosságú a modern technológiákban is, például a mikrohullámú antennákban és a rádióhullámok terjedésében.
\begin{figure}[H]
    \centering
    \includegraphics[width=0.5\textwidth]{imgs/8-tetel/diffrakcio.png}
    \caption{Diffrakció egy résen}
\end{figure}

Ha például egy vékony rést helyezünk el egy sík hullám útjába, akkor a hullám a résen áthaladva elhajlik, és új hullámfrontok jönnek létre a rés mögött. Ezek az új hullámfrontok interferálnak egymással, és diffrakciós mintázatot hoznak létre a térben.

\subsubsection*{Diffrakció véges nyíláson}
Ha egy véges szélességű $a$ nyíláson halad át egy sík hullám, akkor a nyílás minden pontja új hullámforrásként viselkedik, és szférikus hullámok indulnak ki ezekből a pontokból.
\begin{figure}[H]
    \centering
    \includegraphics[width=0.5\textwidth]{imgs/8-tetel/diffrakcio_veges_nyilason.png}
    \caption{Diffrakció véges nyíláson}
\end{figure}
Ha ekkor az $y$ tengely mentén (ami párhuzamos a nyílással) veszünk egy kis $dy$ elemet, akkor ennek a kis elemnek az amplitúdója egy távoli $P$ pontban:
\begin{equation}
    du_p = \frac{u_0}{r} e^{i(\omega t - k r)} dy
\end{equation}
A teljes amplitúdó pedig az $-a/2$ és $a/2$ között integrálva:
\begin{equation}
    u_p = \frac{u_0}{r} e^{i(\omega t - k r)} \int_{-a/2}^{a/2} e^{i k y \sin \theta} dy
\end{equation}
Ahol $r$ a távolság a nyílástól a $P$ pontig, és $\theta$ a $P$ pont és a középvonal által bezárt szög. Az integrál kiszámolásával az intenzitás a következőképpen alakul:
\begin{equation}
    I_p = I(0) \text{sinc}^2 \left( \frac{ka \sin \theta}{2} \right)
\end{equation}
Ahol $\text{sinc}(x) = \frac{\sin x}{x}$ és $I(0)$ az intenzitás a középvonalon ($\theta = 0$). Ez az eredmény azt mutatja, hogy a véges nyíláson áthaladó hullám diffrakciós mintázatot hoz létre, amelynek főmaximuma a középvonalon van, és oldalsó maximumok és minimumok is megjelennek a térben.
\begin{figure}[H]
    \centering
    \includegraphics[width=0.5\textwidth]{imgs/8-tetel/diffrakcio_intenzitas.png}
    \caption{Diffrakciós intenzitás véges nyíláson}
\end{figure}
Az első minimum helye a következőképpen adódik:
\begin{equation}
    \frac{ka \sin \theta}{2} = \pi \qquad \Rightarrow \qquad \sin \theta = \frac{\lambda}{a}
\end{equation}
Ahol $\lambda$ a hullámhossz. Ez azt jelenti, hogy a diffrakciós mintázat szélessége fordítottan arányos a nyílás szélességével: minél kisebb a nyílás, annál szélesebb a diffrakciós mintázat.

A diffrakciók egyik speciális esete a \textbf{Fraunhofer-diffrakció}, amely akkor fordul elő, amikor a hullámforrás és a megfigyelő távol vannak a diffrakciós akadálytól. Ebben az esetben a hullámfrontok közel párhuzamosak, és a diffrakciós mintázat egyszerűbben elemezhető. A Fraunhofer-diffrakciót gyakran használják optikai rendszerek tervezésében és elemzésében, például lencsék és távcsövek esetén.

\subsection{Retardált potenciálok. Dipólsugárzás}
\subsubsection{Retardált potenciálok}
Statikus esetben a D'Alambert egyenletek a Poisson-egyenletekre egyszerűsödnek:
\begin{equation}
    \boldsymbol{\nabla}^2 \phi = - \frac{\rho}{\epsilon_0} \qquad \text{és} \qquad \boldsymbol{\nabla}^2 \bold{A} = - \mu_0 \bold{J}
\end{equation}
A jól ismert megoldásokkal:
\begin{equation}
    V(\bold{r}) = \frac{1}{4 \pi \epsilon_0} \int \frac{\rho(\bold{r}')}{\rcurs} dV' \qquad \text{és} \qquad \bold{A}(\bold{r}) = \frac{\mu_0}{4 \pi} \int \frac{\bold{J}(\bold{r}')}{\rcurs} dV'
\end{equation}
Ahol $\rcurs = |\bold{r} - \bold{r}'|$. Azt tudjuk, hogy az elektromágneses "hírek" fénysebességgel terjednek. A nem statikus esetben tehát nem az számít, hogy egy
egy adott pillanatban milyen állapotban van a forrás, hanem, hogy egy korábbi állapotban milyen volt. Ezt a $t_r$ időt nevezzük \textbf{retardált időnek}, és a következőképpen definiáljuk:
\begin{equation}
    \boxed{t_r \equiv t - \frac{\rcurs}{c}}
\end{equation}
Ahol $t$ a megfigyelő időpontja, és $\frac{\rcurs}{c}$ az az idő, ami alatt a fény eljut a forrástól a megfigyelőhöz. Ekkor a potenciálok általános megoldása a következőképpen alakul:
\begin{equation}
    \boxed{V(\bold{r}, t) = \frac{1}{4 \pi \epsilon_0} \int \frac{\rho(\bold{r}', t_r)}{\rcurs} dV' \qquad \text{és} \qquad \bold{A}(\bold{r}, t) = \frac{\mu_0}{4 \pi} \int \frac{\bold{J}(\bold{r}', t_r)}{\rcurs} dV'}
\end{equation}
Ezeket nevezzük \textbf{retardált potenciáloknak}.
Itt a $\rho(\bold{r}', t_r)$ a töltéssűrűség a $\bold{r}'$ helyen és a retardált időben. 

\subsubsection*{Bizonyítás:}
De ez valóban igaz? Eddig nem bizonyítottunk semmit se, csak logikusan levezettük abból az állításból, hogy az információ fénysebességgel terjed. Ahhoz, hogy bizonyítsuk,
meg kell mutatni, hogy kielégítik az inhomogén hullámegyenleteket és kielégítik a Lorenz-feltételt. Ez azért fontos, mert ez nem egy egyértelmű dolog,
például ha ezt a logikát a terekre alkalmazzuk, akkor ROSSZ választ fogunk kapni:
\begin{equation}
    \bold{E}(\bold{r}, t) \neq \frac{1}{4 \pi \epsilon_0} \int \frac{\rho(\bold{r}', t_r)}{\rcurs^2} \hat{\rcurs} dV' \qquad \text{és} \qquad \bold{B}(\bold{r}, t) \neq \frac{\mu_0}{4 \pi} \int \frac{\bold{J}(\bold{r}', t_r)}{\rcurs^2} \hat{\rcurs} \times dV'
\end{equation}
Ezért meg kell néznünk, hogy mi a retardált potenciál alakulását a Nabla-operátorral. A retardált potenciál két helyen függ $\bold{r}$-től: az egyik a $\rcurs$ a nevezőben, a másik pedig a $t_r = t - \frac{\rcurs}{c}$ a forrásfüggvényekben. 
Tehát a potenciál gradiense a következőképpen alakul:
\begin{equation}
    \boldsymbol{\nabla} \cdot V(\bold{r}, t) = \frac{1}{4 \pi \epsilon_0} \int \left[(\boldsymbol{\nabla} \rho) \frac{1}{\rcurs} + \rho \boldsymbol{\nabla} (\frac{1}{\rcurs})\right] dV'
\end{equation} 
Ahol az első tagban a $\boldsymbol{\nabla} \rho$ a forrásfüggvényekre vonatkozik, és ezt a láncszabály segítségével tudjuk kiszámolni:
\begin{equation}
    \boldsymbol{\nabla} \rho(\bold{r}', t_r) = \frac{\partial \rho}{\partial t_r} \boldsymbol{\nabla} t_r = - \frac{1}{c} \frac{\partial \rho}{\partial t_r} \boldsymbol{\nabla} \rcurs
\end{equation}
Vezessünk pár extra jelölést: $\boldsymbol{\nabla} \rcurs = \hrcurs$ és $\boldsymbol{\nabla} (\frac{1}{\rcurs}) = - \frac{\hrcurs}{\rcurs^2}$. Ekkor a potenciál divergenciája a következőképpen alakul:
\begin{equation}
    \boldsymbol{\nabla} V = \frac{1}{4 \pi \epsilon_0} \int \left[- \frac{1}{c} \frac{\partial \rho}{\partial t_r} \frac{\hrcurs}{\rcurs} - \rho \frac{\hrcurs}{\rcurs^2}\right] dV'
\end{equation}
Most vennünk kell az eredmény divergenciáját (hogy megkapjuk $\boldsymbol{\nabla}^2 V$-t):
\begin{equation}
    \boldsymbol{\nabla}^2 V = \frac{1}{4 \pi \epsilon_0} \int \left[- \frac{1}{c} \boldsymbol{\nabla} \left(\frac{\partial \rho}{\partial t_r} \frac{\hrcurs}{\rcurs}\right) - \boldsymbol{\nabla} \left(\rho \frac{\hrcurs}{\rcurs^2}\right)\right] dV'
\end{equation}
Jelöljük a retardált időderiváltat a Newton jelöléssel ($\frac{\partial \rho}{\partial t_r} = \dot{\rho}$), tehát:
\begin{equation}
    \begin{aligned}
        \boldsymbol{\nabla}^2 V = \frac{1}{4 \pi \epsilon_0} \int &- \frac{1}{c} \left[\frac{\brcurs}{\rcurs} \cdot (\boldsymbol{\nabla} \dot{\rho}) + \dot{\rho} \boldsymbol{\nabla} \cdot \frac{\brcurs}{\rcurs}\right] \\
        &- \left[\frac{\brcurs}{\rcurs^2} \cdot (\boldsymbol{\nabla} \rho) + \rho \boldsymbol{\nabla} \cdot \frac{\brcurs}{\rcurs^2}\right] dV'
    \end{aligned}
\end{equation}
És a láncszabály segítségével a következőképpen írhatjuk át a $\boldsymbol{\nabla} \dot{\rho}$ és $\boldsymbol{\nabla} \cdot \frac{\hrcurs}{\rcurs}$ kifejezéseket:
\begin{equation}
    \boldsymbol{\nabla} \dot{\rho} = - \frac{1}{c} \ddot{\rho} \hrcurs \qquad \text{és} \qquad \boldsymbol{\nabla} \cdot \frac{\hrcurs}{\rcurs} = \frac{1}{\rcurs^2}
\end{equation}
És:
\begin{equation}
    \boldsymbol{\nabla} \cdot (\frac{\hrcurs}{\rcurs^2}) = 4 \pi \delta^3(\brcurs)
\end{equation}
Ahol a $\delta^3(\brcurs)$ a háromdimenziós Dirac-delta függvény. Ezt beírva az előző egyenletbe, a következőképpen alakul $\boldsymbol{\nabla}^2 V$:
\begin{equation}
    \boldsymbol{\nabla}^2 V = \frac{1}{4 \pi \epsilon_0} \int \left[\frac{1}{c^2} \frac{\ddot{\rho}}{\rcurs} - 4 \pi \rho \delta^3(\brcurs)\right] dV' = \frac{1}{c^2} \frac{\partial^2 V}{\partial t^2} - \frac{1}{\epsilon_0}\rho(\bold{r}, t)
\end{equation}
Ez pontosan az inhomogén hullámegyenlet a potenciálra nézve. Érdekesség, hogy ez a Bizonyítás érvényes a fejlett potenciálra is:
\begin{equation}
    V_a(\bold{r}, t) = \frac{1}{4 \pi \epsilon_0} \int \frac{\rho(\bold{r}', t_a)}{\rcurs} dV' \qquad \text{ahol} \qquad t_a = t + \frac{\rcurs}{c}
\end{equation}
Viszont ennek nincs fizikai jelentése, hiszen ez azt jelentené, hogy a jövőbeli információk hatnak a jelenre, ami sérti a kauzalitás elvét.

Most nézzük meg a Lorenz-feltételt:
\begin{equation}
    \boldsymbol{\nabla} \cdot \bold{A} + \mu_0 \epsilon_0 \frac{\partial V}{\partial t} = 0
\end{equation}
A potenciálok behelyettesítése után a következőképpen alakul:
\begin{equation}
    \boldsymbol{\nabla} \cdot \bold{A} = \frac{\mu_0}{4 \pi} \int \left[(\boldsymbol{\nabla} \cdot \bold{J}) \frac{1}{\rcurs} + \bold{J} \cdot \boldsymbol{\nabla} (\frac{1}{\rcurs})\right] dV'
\end{equation}
Ahol az első tagban a $\boldsymbol{\nabla} \cdot \bold{J}$ a forrásfüggvényekre vonatkozik, és ezt a láncszabály segítségével tudjuk kiszámolni:
\begin{equation}
    \boldsymbol{\nabla} \cdot \bold{J}(\bold{r}', t_r) = \frac{\partial \bold{J}}{\partial t_r} \cdot \boldsymbol{\nabla} t_r = - \frac{1}{c} \frac{\partial \bold{J}}{\partial t_r} \cdot \boldsymbol{\nabla} \rcurs
\end{equation}
Ezt beírva az előző egyenletbe, és a már ismert jelöléseket használva ($\boldsymbol{\nabla} \rcurs = \hrcurs$ és $\boldsymbol{\nabla} (\frac{1}{\rcurs}) = - \frac{\hrcurs}{\rcurs^2}$), a potenciál divergenciája a következőképpen alakul:
\begin{equation}
    \boldsymbol{\nabla} \cdot \bold{A} = \frac{\mu_0}{4 \pi} \int \left[- \frac{1}{c} \frac{\partial \bold{J}}{\partial t_r} \cdot \frac{\hrcurs}{\rcurs} - \bold{J} \cdot \frac{\hrcurs}{\rcurs^2}\right] dV'
\end{equation}
Most vegyük a potenciál időderiváltját:
\begin{equation}
    \frac{\partial V}{\partial t} = \frac{1}{4 \pi \epsilon_0} \int \frac{\partial \rho(\bold{r}', t_r)}{\partial t} \frac{1}{\rcurs} dV' = \frac{1}{4 \pi \epsilon_0} \int \frac{\partial \rho(\bold{r}', t_r)}{\partial t_r} \frac{\partial t_r}{\partial t} \frac{1}{\rcurs} dV' = \frac{1}{4 \pi \epsilon_0} \int \frac{\dot{\rho}}{\rcurs} dV'
\end{equation}
Most összeadva a két eredményt:
\begin{equation}
    \boldsymbol{\nabla} \cdot \bold{A} + \mu_0 \epsilon_0 \frac{\partial V}{\partial t} = \frac{\mu_0}{4 \pi} \int \left[- \frac{1}{c} \frac{\partial \bold{J}}{\partial t_r} \cdot \frac{\hrcurs}{\rcurs} - \bold{J} \cdot \frac{\hrcurs}{\rcurs^2} + \frac{1}{c} \frac{\dot{\rho}}{\rcurs}\right] dV'
\end{equation}
Most alkalmazzuk a töltésmegmaradás törvényét a forrásfüggvényekre:
\begin{equation}
    \boldsymbol{\nabla}' \cdot \bold{J}(\bold{r}', t_r) + \frac{\partial \rho(\bold{r}', t_r)}{\partial t_r} = 0
\end{equation}
Ezt beírva az előző egyenletbe, és alkalmazva a láncszabályt a $\boldsymbol{\nabla}' \cdot \bold{J}$ kifejezésre:
\begin{equation}
    \boldsymbol{\nabla}' \cdot \bold{J}(\bold{r}', t_r) = \frac{\partial \bold{J}}{\partial t_r} \cdot \boldsymbol{\nabla}' t_r = - \frac{1}{c} \frac{\partial \bold{J}}{\partial t_r} \cdot \boldsymbol{\nabla}' \rcurs = - \frac{1}{c} \frac{\partial \bold{J}}{\partial t_r} \cdot (-\hrcurs) = \frac{1}{c} \frac{\partial \bold{J}}{\partial t_r} \cdot \hrcurs
\end{equation}
Ezt beírva az előző egyenletbe, a következőképpen alakul:
\begin{equation}
    \boldsymbol{\nabla} \cdot \bold{A} + \mu_0 \epsilon_0 \frac{\partial V}{\partial t} = \frac{\mu_0}{4 \pi} \int \left[- \boldsymbol{\nabla}' \cdot \bold{J} \frac{1}{\rcurs} - \bold{J} \cdot \frac{\hrcurs}{\rcurs^2} + \frac{1}{c} \frac{\dot{\rho}}{\rcurs}\right] dV'
\end{equation}
Most alkalmazzuk a töltésmegmaradás törvényét az utolsó két tagra is:
\begin{equation}
    \boldsymbol{\nabla} \cdot \bold{A} + \mu_0 \epsilon_0 \frac{\partial V}{\partial t} = \frac{\mu_0}{4 \pi} \int \left[- \boldsymbol{\nabla}' \cdot \bold{J} \frac{1}{\rcurs} + \boldsymbol{\nabla}' \cdot \bold{J} \frac{1}{\rcurs}\right] dV' = 0
\end{equation}
Tehát a retardált potenciálok kielégítik a Lorenz-feltételt is.

\subsubsection{Dipólsugárzás}
Mi a sugárzás? Amikor a töltések gyorsulnak, az elektromágneses terük képes energiát szállítani a végtelenbe. Ezt a jelenséget nevezzük \textbf{sugárzásnak}.
Tegyük fel, hogy a forrás nagyjából lokalizált az eredet környékén. Ekkor meg akarom határozni az energiát ami $t_0$ időben sugárzódik ki. Vegyünk egy nagy gömbfelületet,
amelynek a sugara $r$, és amelynek a középpontja az eredetben van. Az ezen a felületen átmenő teljesítmény a Poynting-vektor integráljával adható meg:
\begin{equation}
    P(t_0) = \oint \bold{S}(\bold{r}, t_0) \cdot d\bold{a} = \oint \frac{1}{\mu_0} \bold{E} \times \bold{B} \cdot d\bold{a}
\end{equation} 
Mivel az információ fénysebességgel halad, ezért ennek az energiának kicsit korábban kellett elindulnia $t_0 = t - \frac{r}{c}$ időben. Tehát a kisugárzott teljesítmény:
\begin{equation}
    P_{rad}(t_0) = \lim_{r \to \infty} P(r, t_0 + \frac{r}{c})
\end{equation}
Ahol $t_0$ egy konstans maradt. Ez az energia (egységnyi idő alatt) távozik és sose jön vissza. 

Mivel a gömbfelület területe $4 \pi r^2$, ezért a Poynting-vektor nagy $r$-nél maximum $1/r^2$-el tud csökkenni. A Coulomb törvény szerint az elektrosztatikus tér is $1/r^2$-el csökken
és a Biot-Savart törvény szerint a mágneses tér is $1/r^2$-el (vagy gyorsabban) csökken. Tehát a Poynting-vektor $S \approx 1/r^4$-el (vagy gyorsabban) csökken statikus konfigurációkban.
Ezért statikus konfigurációkban nincs sugárzás. Ahhoz, hogy legyen sugárzás, a tereknek lassabban kell csökkenniük, mint $1/r^2$. De  Jeremenko egyenletek segítségével megmutatható, hogy az
időfüggő tereknek vannak olyan komponensei, amik $1/r$-el csökkennek. Ezek a komponensek felelősek az elektromágneses sugárzásért.

A sugárzás tanulmányozása tehát arról szól, hogy kiszedjük $\bold{E}$ és $\bold{B}$ terek $1/r$-el csökkenő komponenseit (nagy $r$-nél), és ezeket hozzákapcsoljuk a Poynting-vektor $1/r^2$-el csökkenő komponenséhez és kiintegáljuk egy nagy gömbfelületen $r \to \infty$ határértékben.

\subsubsection*{Elektromos dipólsugárzás}
Képzeljünk el egy rrendszert amiben van két fémgömb amik egy vékony vezetékkel össze vannak kötve $d$ távolságban. $t$ időben a felső gömb töltése $+q(t)$, az alsó gömb töltése pedig $-q(t)$.
Most tegyük fel, hogy le-fel mozgatjuk a töltéseket $\omega$ szögfrekvenciával:
\begin{equation}
    q(t) = q_0 \cos(\omega t)
\end{equation}
Ekkor a dipólmomentum a következőképpen alakul:
\begin{equation}
    \bold{p}(t) = q(t) d \hat{z} = p_0 \cos(\omega t) \hat{\bold{z}} \qquad \text{ahol} \qquad p_0 \equiv q_0 d
\end{equation}
Ahol $p_0$ a dipólmomentum maximum értéke. 

A retardált potenciál:
\begin{equation}
    V(\bold{r}, t) = \frac{1}{4 \pi \epsilon_0} \left\{\frac{q_0 \cos \left[\omega(t - \frac{\rcurs_+}{c} )\right]}{\rcurs_+} - \frac{q_0 \cos \left[\omega(t - \frac{\rcurs_-}{c} )\right]}{\rcurs_-}\right\}
\end{equation}
Ahol $\rcurs_\pm$ a cosinus-tétel alapján a következőképpen alakul:
\begin{equation}
    \rcurs_\pm = \sqrt{r^2 + \left(\frac{d}{2}\right)^2 \mp r d \cos \theta} \approx r \mp \frac{d}{2} \cos \theta
\end{equation}
\begin{figure}[H]
    \centering
    \includegraphics[width=0.5\textwidth]{imgs/8-tetel/elektromos_dipol_szogek.png}
    \caption{Elektromos dipólsugárzás geometriája}
\end{figure}

Ahhoz, hogy egy fizikai dipólust egy ideális dipólussá alakítsunk, a szeparációs távolságot nagyon kicsire kell vennünk:
\begin{equation}
    \text{\textbf{közelítés 1:}} \qquad d \ll r
\end{equation}
Természetesen ha $d = 0$, akkor nincs potenciál, amit akarunk az egy expanzió ami a $d$-ben elsőrendű, tehát:
\begin{equation}
    \rcurs_\pm \approx r \left(1 \mp \frac{d}{2r} \cos \theta \right)
\end{equation}
Vagyis:
\begin{equation}
    \frac{1}{\rcurs_\pm} \approx \frac{1}{r} \left(1 \pm \frac{d}{2r} \cos \theta \right)
\end{equation}
Tehát a $cos$ függvény:
\begin{equation}
    \begin{aligned}
        \cos &\left[\omega (t - \frac{\rcurs_\pm}{c})\right] \approx \\
        &\approx \cos \left[\omega \left(t - \frac{r}{c}\right) \pm \frac{\omega d}{2c} \cos \theta \right] = \\
        &= \cos \left[\omega (t - \frac{r}{c})\right] \cos \left(\frac{\omega d}{2 c} \cos \theta \right) \mp \sin \left[\omega (t - \frac{r}{c})\right] \sin \left(\frac{\omega d}{2 c} \cos \theta \right)
    \end{aligned}
\end{equation}
Most tegyük fel, hogy a dipólus nagyon kicsi a hullámhosszhoz képest:
\begin{equation}
    \text{\textbf{közelítés 2:}} \qquad d \ll \lambda = \frac{2 \pi c}{\omega} \implies d \ll \frac{c}{\omega}
\end{equation}
Ekkor a koszinusz függvényt tovább tudjuk közelíteni:
\begin{equation}
    \cos \left[\omega (t - \frac{\rcurs_\pm}{c})\right] \approx \cos \left[\omega (t - \frac{r}{c})\right] \mp \frac{\omega d}{2 c} \cos \theta \sin \left[\omega (t - \frac{r}{c})\right]
\end{equation}
Ezeket visszahelyettesítve a potenciálba, és csak az elsőrendű tagokat megtartva, a következőképpen alakul a potenciál:
\begin{equation}
    V(\bold{r}, \theta, t) = \frac{p_0 \cos \theta}{4 \pi \epsilon_0 r^2} \left\{- \frac{\omega}{c} \sin \left[\omega (t - \frac{r}{c})\right] + \frac{1}{r}\cos\left[\omega (t - \frac{r}{c})\right] \right\} 
\end{equation}
Statikus esetben ($\omega \to 0$) természetesen visszakapjuk a dipólus potenciált:
\begin{equation}
    V_{stat}(\bold{r}, \theta) = \frac{p_0 \cos \theta}{4 \pi \epsilon_0 r^2}
\end{equation}
Ez viszont nem az a régió ami minket érdekel. Mi arra vagyunk kíváncsiak, hogy milyen terek élik túl a nagy távolságot a forrástól, amit \textbf{sugárzási zónának} nevezünk:
\begin{equation}
    \text{\textbf{közelítés 3:}} \qquad r \gg \lambda = \frac{2 \pi c}{\omega} \implies r \gg \frac{c}{\omega}
\end{equation}
Ebben a régióban a potenciál a következő alakra egyszerűsödik:
\begin{equation}
    \boxed{V_{rad}(\bold{r}, \theta, t) = - \frac{p_0 \omega \cos \theta}{4 \pi \epsilon_0 c r} \sin \left[\omega (t - \frac{r}{c})\right]}
\end{equation}
A vektorpotenciál pedig meghatározza a vezetékben folyó áramot:
\begin{equation}
    \bold{A}(\bold{r}, t) = \frac{\mu_0}{4 \pi} \int_{-d/2}^{d/2} \frac{-q_0 \omega \sin \left[\omega (t - \frac{\rcurs}{c})\right] \hat{\bold{z}}}{\rcurs} dz
\end{equation}
Tehát:
\begin{equation}
    \bold{I}(\bold{r}, t) = \frac{d q}{d t} \hat{\bold{z}} = - q_0 \omega \sin \left[\omega (t - \frac{\rcurs}{c})\right] \hat{\bold{z}}
\end{equation}
Mivel az integrál bevezet egy faktort $d$-nek, ezért első rendben közelíthetünk azzal, hogy az integrált kicseréljük az értékével:
\begin{equation}
    \boxed{\bold{A}_{rad}(\bold{r}, \theta, t) = - \frac{\mu_0 p_0 \omega}{4 \pi r} \sin \left[\omega (t - \frac{r}{c})\right] \hat{\bold{z}}}
\end{equation}
Ebből a potenciálból pedig könnyen maghatárrozhatóak a terek a sugárzási zónában:
\begin{equation}
    \boxed{\bold{E}_{rad}(\bold{r}, \theta, t) = - \boldsymbol{\nabla} V - \frac{\partial \bold{A}}{\partial t} = - \frac{\mu_0 p_0 \omega^2}{4 \pi} \left(\frac{\sin \theta}{r}\right) \cos \left[\omega (t - \frac{r}{c})\right] \hat{\boldsymbol{\theta}}}
\end{equation}
\begin{equation}
    \boxed{\bold{B}_{rad}(\bold{r}, \theta, t) = \boldsymbol{\nabla} \times \bold{A} = - \frac{\mu_0 p_0 \omega^2}{4 \pi c} \left(\frac{\sin \theta}{r}\right) \cos \left[\omega (t - \frac{r}{c})\right] \hat{\boldsymbol{\phi}}}
\end{equation}
Ahol a $\hat{\boldsymbol{\theta}}$ és $\hat{\boldsymbol{\phi}}$ a gömbi koordinátarendszer egységvektorai. Ezek a terek $1/r$-el csökkennek, és ezért képesek energiát szállítani a végtelenbe.
Ezek az egyenletek monokromatikus hullámokat eprezentálnak $\omega$ frekvenciával, amelyek a dipólus tengelyére merőlegesen sugároznak. Az amplitúdóik hányadosa pedig $\frac{E_0}{B_0} = c$.
A terek fázisban vannak egymással, merőlegesek egymásra és transzverzálisak. Tehát pontosan úgy viselkednek ahogy gondoltuk elektromágneses terek esetében vákuumban. Igazából, a valóságban
ezek gömbhullámok lesznek, de elég távol jó közelítéssel síkhullámokként viselkednek.

Az energia ami kisugárzik a dipólusból a Poynting-vektor segítségével határozható meg:
\begin{equation}
    \bold{S}_{rad}(\bold{r}, \theta, t) = \frac{1}{\mu_0} \bold{E}_{rad} \times \bold{B}_{rad} = \frac{\mu_0 p_0^2 \omega^4}{16 \pi^2 c} \left(\frac{\sin^2 \theta}{r^2}\right) \cos^2 \left[\omega (t - \frac{r}{c})\right] \hat{\boldsymbol{r}}
\end{equation}
Az intenzitrás pedig a Poynting-vektor időátlagolt értéke egy adott irányban:
\begin{equation}
    \langle S_{rad}(\bold{r}, \theta) \rangle = \frac{1}{T} \int_0^T \bold{S}_{rad}(\bold{r}, \theta, t) dt = \frac{\mu_0 p_0^2 \omega^4}{32 \pi^2 c} \left(\frac{\sin^2 \theta}{r^2}\right) \hat{\boldsymbol{r}}
\end{equation}
Ahol $T = \frac{2 \pi}{\omega}$ a hullám periódusa. Vegyük észre, hogy nincs sugárzás a dipólus tengelyének irányában ($\theta = 0$ vagy $\pi$), hiszen ott nincs gyorsuló töltés komponens.
Ez azt jelenti, hogy az intenzitás profil egy toroid alakú minta lesz a dipólus körül, az egyenlítő síkjában maximum értékkel.
\begin{figure}[H]
    \centering
    \includegraphics[width=0.5\textwidth]{imgs/8-tetel/dipol_sugarzas_intenzitas.png}
    \caption{Dipólus sugárzás intenzitás profilja}
\end{figure}
A teljes kisugárzott teljesítmény pedig a következőképpen alakul:
\begin{equation}
    \langle P_{rad} \rangle = \oint \langle \bold{S}_{rad}(\bold{r}, \theta) \rangle \cdot d\bold{A} = \frac{\mu_0 p_0^2 \omega^4}{12 \pi c}
\end{equation}

\subsubsection*{Mágneses Dipólsugárzás}
Most vegyünk egy vezető hurkot aminek $b$ a sugara, és $I(t)$ a váltakozó áram benne:
\begin{equation}
    I(t) = I_0 \cos(\omega t)
\end{equation}
Ez pedig az oszcilláló mágneses dipólusmomentum modellje:
\begin{equation}
    \bold{m}(t) = I(t) \pi b^2 \hat{z} = m_0 \cos(\omega t) \hat{z} \qquad \text{ahol} \qquad m_0 \equiv I_0 \pi b^2
\end{equation}
Ahol $m_0$ a mágneses dipólmomentum maximum értéke.
\begin{figure}
    \centering
    \includegraphics[width=0.5\textwidth]{imgs/8-tetel/magnetikus_dipol_szogek.png}
    \caption{Mágneses dipólsugárzás geometriája}
\end{figure}
A hurok nem változik ezért a skalár potenciál nulla. A retardált vektorpotenciál pedig a következőképpen alakul:
\begin{equation}
    \bold{A}(\bold{r}, t) = \frac{\mu_0}{4 \pi} \int \frac{\bold{J}(\bold{r}', t_r)}{\rcurs} dV' = \frac{\mu_0}{4 \pi} \int \frac{I_0  cos \left[\omega (t - \frac{\rcurs}{c})\right]}{\rcurs} d\bold{I}'
\end{equation}
Egy pontra közvetlenül az $x$ tengely fölött. $\bold{A}$-nak $y$ irányba kell mutatnia mivel az $x$ komponens szimmetria miatt kiesik, és a $z$ komponens is kiesik mivel a $\bold{I}'$ és $\hat{\rcurs}$ vektorok merőlegesek egymásra.
Tehát a vektorpotenciál a sugárzási zónában a következőképpen alakul:
\begin{equation}
    \bold{A}_{rad}(\bold{r}, \theta, t) = \frac{\mu_0 I_0 b}{4 \pi} \hat{\bold{y}} \int_0^{2 \pi} \frac{\cos \left[\omega (t - \frac{\rcurs}{c})\right]}{\rcurs} \cos \phi' d\phi'
\end{equation}
A $cos \phi'$ azért kell, hogy kiválassza az $y$ komponensét $d\bold{I}'$-nek. A Koszinusz tétel alapján a következőképpen alakul $\rcurs$:
\begin{equation}
    \rcurs = \sqrt{r^2 + b^2 - 2 r b \cos \psi}
\end{equation}
Ahol $\psi$ az a szög ami az $\bold{r}$ és $\bold{b}$ vektorok között van:
\begin{equation}
    \bold{r} = r \sin \theta \hat{\bold{x}} + r \cos \theta \hat{\bold{z}} \qquad \text{és} \qquad \bold{b} = b \cos \phi' \hat{\bold{x}} + b \sin \phi' \hat{\bold{y}}
\end{equation}
Tehát $\psi = \bold{r} \cdot \bold{b} = r b \sin \theta \cos \phi'$, ezért:
\begin{equation}
    \rcurs = \sqrt{r^2 + b^2 - 2 r b \sin \theta \cos \phi'}
\end{equation}
Most megint be kell vezetnünk néhány közelítést, ahhoz, hogy ideális dipólusokkal tudjunk számolni. Először is a hurkot nagyon kicsire vesszük a távolsághoz képest:
\begin{equation}
    \text{\textbf{közelítés 1:}} \qquad b \ll r
\end{equation}
Ekkor:
\begin{equation}
    \rcurs \approx r \left(1 - \frac{b}{r} \sin \theta \cos \phi' \right)
\end{equation}
Vagyis:
\begin{equation}
    \frac{1}{\rcurs} \approx \frac{1}{r} \left(1 + \frac{b}{r} \sin \theta \cos \phi' \right)
\end{equation}
És a koszinusz függvény:
\begin{equation}
    \begin{aligned}
        \cos &\left[\omega (t - \frac{\rcurs}{c})\right] \approx \\
        &\approx \cos \left[\omega \left(t - \frac{r}{c}\right) + \frac{\omega b}{c} \sin \theta \cos \phi' \right] = \\
        &= \cos \left[\omega (t - \frac{r}{c})\right] \cos \left(\frac{\omega b}{c} \sin \theta \cos \phi'\right) = \\
        &= \cos \left[\omega (t - \frac{r}{c})\right] \cos \left(\frac{\omega b}{c} \sin \theta \cos \phi' \right) - \sin \left[\omega (t - \frac{r}{c})\right] \sin \left(\frac{\omega b}{c} \sin \theta \cos \phi' \right)
    \end{aligned}
\end{equation}
Mint az előbb, most is feltesszük, hogy a dipólus kicsi a sugárzás hullámhosszához képest:
\begin{equation}
    \text{\textbf{közelítés 2:}} \qquad b \ll \frac{c}{\omega}
\end{equation}
Ekkor a koszinusz függvényt tovább tudjuk közelíteni:
\begin{equation}
    \cos \left[\omega (t - \frac{\rcurs}{c})\right] \approx \cos \left[\omega (t - \frac{r}{c})\right] - \frac{\omega b}{c} \sin \theta \cos \phi' \sin \left[\omega (t - \frac{r}{c})\right]
\end{equation}
Ha behelyettesítjük a $1/\rcurs$ és a koszinusz közelítéseket a vektorpotenciálba, és csak az elsőrendű tagokat tartjuk meg, akkor a következőképpen alakul a vektorpotenciál:
\begin{equation}
    \begin{aligned}
        \bold{A}(\bold{r}, \theta, t) \approx \frac{\mu_0 I_0 b}{4 \pi r} \hat{\bold{y}} \int_0^{2 \pi} &\left\{\cos \left[\omega (t - \frac{r}{c}) \right] + b \sin \theta \cos \phi' \times \right. \\
        &\left. \times \left(\frac{1}{r} \cos \left[\omega (t - \frac{r}{c})\right]  - \frac{\omega}{c} \sin \left[\omega (t - \frac{r}{c})\right]\right) \right\} \cos \phi' d\phi'
    \end{aligned}
\end{equation}
Az első tag untegrálja nulla, mivel az integrál egy teljes perióduson megy végig:
\begin{equation}
    \int_0^{2 \pi} \cos \phi' d\phi' = 0
\end{equation}
A második tagban $cos^2 \phi'$ szerepel, aminek az integrálja:
\begin{equation}
    \int_0^{2 \pi} \cos^2 \phi' d\phi' = \pi
\end{equation}
Ezt beírva a vektorpotenciálba, figyelembe véve, hogy általában a $\bold{A}$ a $\hat{\boldsymbol{\phi}}$ irányba mutat a sugárzási zónában, a következőképpen alakul a vektorpotenciál:
\begin{equation}
    \bold{A}_{rad}(\bold{r}, \theta, t) = \frac{\mu_0 m_0}{4 \pi} \left(\frac{sin \theta}{r}\right) \left\{\frac{1}{r} \cos \left[\omega (t - \frac{r}{c})\right] - \frac{\omega}{c} \sin \left[\omega (t - \frac{r}{c})\right] \right\} \hat{\boldsymbol{\phi}}
\end{equation}
Ez a tökéletes oszcilláló mágneses dipólus vektorpotenciálja a sugárzási zónában. Statikus esetben ($\omega = 0$) visszakapjuk a mágneses dipólus vektorpotenciált:
\begin{equation}
    \bold{A}_{stat}(\bold{r}, \theta) = \frac{\mu_0 m_0}{4 \pi} \left(\frac{sin \theta}{r^2}\right) \hat{\boldsymbol{\phi}}
\end{equation}
De minket a sugárzási zóna érdekel, tehát még egy közelítést be kell vezetnünk:
\begin{equation}
    \text{\textbf{közelítés 3:}} \qquad r \gg \frac{c}{\omega}
\end{equation}
Tehát a $\bold{A}$ első tagja elhanyagolható a másodikhoz képest, és a vektorpotenciál a következőképpen alakul:
\begin{equation}
    \boxed{\bold{A}_{rad}(\bold{r}, \theta, t) = - \frac{\mu_0 m_0 \omega}{4 \pi c} \left(\frac{sin \theta}{r}\right) \sin \left[\omega (t - \frac{r}{c})\right] \hat{\boldsymbol{\phi}}}
\end{equation}
Ebből a potenciálból pedig könnyen meghatározhatóak a terek a sugárzási zónában:
\begin{equation}
    \boxed{\bold{E}_{rad}(\bold{r}, \theta, t) = - \frac{\partial \bold{A}}{\partial t} = - \frac{\mu_0 m_0 \omega^2}{4 \pi c} \left(\frac{sin \theta}{r}\right) \cos \left[\omega (t - \frac{r}{c})\right] \hat{\boldsymbol{\phi}}}
\end{equation}
\begin{equation}
    \boxed{\bold{B}_{rad}(\bold{r}, \theta, t) = \boldsymbol{\nabla} \times \bold{A} = - \frac{\mu_0 m_0 \omega^2}{4 \pi c^2} \left(\frac{sin \theta}{r}\right) \cos \left[\omega (t - \frac{r}{c})\right] \hat{\boldsymbol{\theta}}}
\end{equation}
Ahol a $\hat{\boldsymbol{\theta}}$ és $\hat{\boldsymbol{\phi}}$ a gömbi koordinátarendszer egységvektorai. Ezek a terek $1/r$-el csökkennek, és ezért képesek energiát szállítani a végtelenbe.
Ezek az egyenletek monokromatikus hullámokat eprezentálnak $\omega$ frekvenciával, amelyek a dipólus tengelyére merőlegesen sugároznak. Az amplitúdóik hányadosa pedig $\frac{E_0}{B_0} = c$.
A terek fázisban vannak egymással, merőlegesek egymásra és transzverzálisak. Tehát pontosan úgy viselkednek ahogy gondoltuk elektromágneses terek esetében vákuumban. Az energia ami kisugárzik a dipólusból a Poynting-vektor segítségével határozható meg:
\begin{equation}
    \bold{S}_{rad}(\bold{r}, \theta, t) = \frac{1}{\mu_0} \bold{E}_{rad} \times \bold{B}_{rad} = \frac{\mu_0 m_0^2 \omega^4}{16 \pi^2 c^3} \left(\frac{\sin^2 \theta}{r^2}\right) \cos^2 \left[\omega (t - \frac{r}{c})\right] \hat{\boldsymbol{r}}
\end{equation}
Az intenzitrás pedig a Poynting-vektor időátlagolt értéke egy adott irányban:
\begin{equation}
    \langle S_{rad}(\bold{r}, \theta) \rangle = \frac{1}{T} \int_0^T \bold{S}_{rad}(\bold{r}, \theta, t) dt = \frac{\mu_0 m_0^2 \omega^4}{32 \pi^2 c^3} \left(\frac{\sin^2 \theta}{r^2}\right) \hat{\boldsymbol{r}}
\end{equation}
Ahol $T = \frac{2 \pi}{\omega}$ a hullám periódusa. Vegyük észre, hogy itt sincs sugárzás a dipólus tengelyének irányában ($\theta = 0$ vagy $\pi$), hiszen ott nincs gyorsuló töltés komponens.
Ez azt jelenti, hogy az intenzitás profil egy toroid alakú minta lesz a dipólus körül, az egyenlítő síkjában maximum értékkel.
A teljes kisugárzott teljesítmény pedig a következőképpen alakul:
\begin{equation}
    \langle P_{rad} \rangle = \oint \langle \bold{S}_{rad}(\bold{r}, \theta) \rangle \cdot d\bold{A} = \frac{\mu_0 m_0^2 \omega^4}{12 \pi c^3}
\end{equation}
Viszont van egy fontos különbség az elektromos és a mágneses dipólus sugárzás között: Egy olyan konfigurációban, ahol a dimenziók hasonlóak, a teljesítmény az elektromos dipólsugárzás esetében sokkal nagyobb lesz, mint a mágneses dipólsugárzás esetében:
\begin{equation}
    \frac{P_{\text{mágneses}}}{P_{\text{elektromos}}} = \left(\frac{m_0}{p_0c}\right)^2
\end{equation}
Ahol $m_0 = \pi b^2 I_0$ a mágneses dipólmomentum maximum értéke, és $p_0 = q_0 d$ az elektromos dipólmomentum maximum értéke. Az áram az elektromos esetben $I_0 = q_0 \omega$ volt, és ha $d = \pi b$-nek vesszük, akkor:
\begin{equation}
    \frac{P_{\text{mágneses}}}{P_{\text{elektromos}}} = \left(\frac{\omega b}{c}\right)^2
\end{equation}
De a $\omega b / c$ pont az a mennyiség amit kicsinek feltételeztünk, és itt négyzetre van emelve. Tehát normál esetben az elektromos tér dominál. A mágneses sugárzást csak akkor lehet érzékelni,
ha nagyon óvatosan úgy tervezzük meg a rendszert, hogy minden elektromos hozzájárulást kiiktassunk.

\subsubsection*{Sugárzás tetszőleges forrásból}
Eddig olyan rendszereket vizsgáltunk ahol a mágneses és az elektomos dipólusok oszcilláltak. De a dipólsugárzást nem túl nehéz általánosítani egy teljesen tetszőleges töltés és áram konfigurációra,
feltéve, hogy a forrás relatíve lokalizált az eredet környékén. 
\begin{figure}[H]
    \centering
    \includegraphics[width=0.5\textwidth]{imgs/8-tetel/tetszoleges_forras.png}
    \caption{Tetszőleges lokalizált forrás geometriája}
\end{figure}
Ekkor a retardált potenciálok a következőképpen alakulnak nagy $r$-nél:
\begin{equation}
    V_{rad}(\bold{r}, t) = \frac{1}{4 \pi \epsilon_0} \int \frac{\rho(\bold{r}', t - \frac{\rcurs}{c})}{\rcurs} dV'
\end{equation}
Ahol:
\begin{equation}
    \rcurs = |\bold{r} - \bold{r}'| = \sqrt{r^2 + r'^2 - 2 \bold{r} \cdot \bold{r}'}
\end{equation}
Ahogy eddig is, itt is bevezetjük a következő közelítést:
\begin{equation}
    \text{\textbf{közelítés 1:}} \qquad r' \ll r
\end{equation}
Itt $r'$ a változó ami szerint integrálunnk, az első közelítés ebben a kontextusban azt jelenti, hogy $r'$ maximum értéke sokkal kisebb, mint $r$. Ekkor:
\begin{equation}
    \rcurs \approx r\left(1 - \frac{\bold{r} \cdot \bold{r'}}{r^2}\right)
\end{equation}
Vagyis:
\begin{equation}
    \frac{1}{\rcurs} \approx \frac{1}{r} \left(1 + \frac{\bold{r} \cdot \bold{r'}}{r^2}\right)
\end{equation}
És a töltéssűrűség:
\begin{equation}
    \rho(\bold{r}', t - \frac{\rcurs}{c}) \approx \rho \left(\bold{r'}, t - \frac{r}{c} + \frac{\hat{\bold{r}} \cdot \bold{r'}}{c}\right)
\end{equation}
Ezt a töltéssűrűséget a retardált idő ($t_0 \equiv t - \frac{r}{c}$) körül Taylor sorba fejthetjük:
\begin{equation}
    \rho(\bold{r}', t - \frac{\rcurs}{c}) \approx \rho(\bold{r}', t_0) + \left(\frac{\hat{\bold{r}} \cdot \bold{r'}}{c} \right) \frac{\partial \rho(\bold{r}', t_0)}{\partial t} + \ldots
\end{equation}
A magasabb deriváljtainakltakat eldobhatjuk, ha a következő közelítés érvényes:
\begin{equation}
    \text{\textbf{közelítés 2:}} \qquad r' \ll \frac{c}{|\ddot{\rho}/\dot{\rho}|}, \quad \frac{c}{|\dddot{\rho}/\dot{\rho}|^{1/2}}, \quad \frac{c}{|\ddddot{\rho}/\dot{\rho}|^{1/3}}, \quad \ldots
\end{equation}
Ahol az időderviáltat megint a Newton jelöléssel jelöljük. Egy oszcilláló rendszerre ezek mind $\sim \frac{c}{\omega}$-ra redukálódnak, tehát visszakapjuk a korábbi közelítést.
Ezeket a közelítéseket visszahelyettesítve a potenciálba, és csak az elsőrendű tagokat megtartva, a következőképpen alakul a potenciál:
\begin{equation}
    V_{rad}(\bold{r}, t) \approx \frac{1}{4 \pi \epsilon_0 r} \left[\int \rho(\bold{r'}, t_0) dV' + \frac{\hat{\bold{r}}}{r} \int \bold{r'} \rho(\bold{r'}, t_0) dV' + \frac{\hat{\bold{r}}}{c} \frac{d}{dt} \int \bold{r'} \rho(\bold{r'}, t_0) dV' \right]
\end{equation}
Az első integrál a teljes töltést $Q$ adja meg $t_0$ időpontban. Mivel a töltés időben megmarad, a másik két integrál az elektromos dipólmomentumot adja meg $t_0$-ban. Tehát:
\begin{equation}
    V_{rad}(\bold{r}, t) \approx \frac{1}{4 \pi \epsilon_0} \left[ \frac{Q}{r} + \frac{\hat{\bold{r}} \cdot \bold{p}(t_0)}{r^2} + \frac{\hat{\bold{r}} \cdot \dot{\bold{p}}(t_0)}{rc} \right]
\end{equation} 
Statikus esetben az első és a második tag a monopól és dipól hozzájárulást adja vissza a multipólus kifejtéséhez a potenciálnak. A harmadik tag pedig nulla lenne.

A vektorpotenciál esetében hasonló közelítéseket alkalmazva, a következőképpen alakul a sugárzási zóna vektorpotenciálja:
\begin{equation}
    \bold{A}_{rad}(\bold{r}, t) = \frac{\mu_0}{4 \pi} \int \frac{\bold{J}(\bold{r'}, t - \frac{\rcurs}{c})}{\rcurs} dV'
\end{equation}
Az elsőrendű közelítés miatt $\rcurs$ felcserélhető $r$-rel az integrálban:
\begin{equation}
    \bold{A}_{rad}(\bold{r}, t) \approx \frac{\mu_0}{4 \pi r} \int \bold{J}(\bold{r'}, t_0) dV'
\end{equation}
A $\bold{J}$ integrálja pedig az áram időderiváltját adja meg a töltésmegmaradás törvénye alapján:
\begin{equation}
    \int \bold{J}(\bold{r'}, t_0) dV' = \frac{d}{dt} \int \bold{r'} \rho(\bold{r'}, t_0) dV' = \dot{\bold{p}}(t_0)
\end{equation}
Tehát a vektorpotenciál a következőképpen alakul:
\begin{equation}
    \bold{A}_{rad}(\bold{r}, t) \approx \frac{\mu_0}{4 \pi} \frac{\dot{\bold{p}}(t_0)}{r}
\end{equation}
Mostmár csak a tereket kell meghatározni, ahol még egy közelítést alkalmazunk:
\begin{equation}
    \text{\textbf{közelítés 3:}} \qquad \frac{1}{r^2} \text{-es tagok elhagyhatóak a $\bold{B}$ és $\bold{E}$ térben}
\end{equation}
Például a Coulomb-tér:
\begin{equation}
    \bold{E} = \frac{1}{4 \pi \epsilon_0} \frac{Q}{r^2} \hat{\bold{r}}
\end{equation}
Nem ad semmit az elektromágneses sugárzáshoz. Igazából a sugárzás teljes egészében azokból a részekből jönnek, amelyeket differenciálunk $t_0$ szerint:
\begin{equation}
    \boldsymbol{\nabla} t_0 = - \frac{1}{c} \boldsymbol{\nabla}r = - \frac{1}{c} \hat{\bold{r}}
\end{equation}
Tehát a potenciál gradiense:
\begin{equation}
    \boldsymbol{\nabla} V_{rad}(\bold{r}, t) = \boldsymbol{\nabla} \left[\frac{1}{4 \pi \epsilon_0} \frac{\hat{\bold{r}} \cdot \dot{\bold{p}}(t_0)}{rc}\right] \approx \frac{1}{4 \pi \epsilon_0} \left[\frac{\hat{\bold{r}} \cdot \ddot{\bold{p}}(t_0)}{rc}\right]\boldsymbol{\nabla} t_0 = - \frac{1}{4 \pi \epsilon_0 c^2} \frac{\left[\hat{\bold{r}} \cdot \ddot{\bold{p}}(t_0)\right]}{r} \hat{\bold{r}}
\end{equation}
Hasonlóan a vektor potenciál rotációja:
\begin{equation}
    \boldsymbol{\nabla} \times \bold{A}_{rad}(\bold{r}, t) \approx \frac{\mu_0}{4 \pi r} \left[\boldsymbol{\nabla} \times \dot{\bold{p}}(t_0) \right] = - \frac{\mu_0}{4 \pi c r} \left[\hat{\bold{r}} \times\dot{\bold{p}}(t_0) \right]
\end{equation}
Miközben a vektor potenciál időderiváltja:
\begin{equation}
    \frac{\partial \bold{A}_{rad}(\bold{r}, t)}{\partial t} \approx \frac{\mu_0}{4 \pi} \frac{\ddot{\bold{p}}(t_0)}{r}
\end{equation}
Ezeket visszahelyettesítve az elektromos és mágneses tér kifejezéseibe, a következőképpen alakulnak a terek a sugárzási zónában:
\begin{equation}
    \boxed{\bold{E}_{rad}(\bold{r}, t) = - \boldsymbol{\nabla} V_{rad} - \frac{\partial \bold{A}_{rad}}{\partial t} = \frac{\mu_0}{4 \pi r} \left[(\hat{\bold{r}} \cdot \ddot{\bold{p}})\hat{\bold{r}} - \ddot{\bold{p}}\right] = - \frac{\mu_0}{4 \pi r} \left[\hat{\bold{r}} \times (\hat{\bold{r}} \times \ddot{\bold{p}})\right]}
\end{equation}
Ahol $\ddot{\bold{p}}$ a $t_0 = t - \frac{r}{c}$ időpontban értékelt második időderivált. A mágneses tér pedig:
\begin{equation}
    \boxed{\bold{B}_{rad}(\bold{r}, t) = \boldsymbol{\nabla} \times \bold{A}_{rad} = - \frac{\mu_0}{4 \pi r c} \left[\hat{\bold{r}} \times \ddot{\bold{p}}\right]}
\end{equation}
Ha gömbi-koordinátákat használunk, és a $z$ tengely irányába helyezzük a dipólusmomentumot, akkor a terek a következőképpen alakulnak:
        \begin{equation}
            \bold{E}_{rad}(\bold{r}, \theta, t) = - \frac{\mu_0 \ddot{p}(t_0)}{4 \pi} \left(\frac{\sin \theta}{r}\right) \hat{\boldsymbol{\theta}}
\end{equation}
\begin{equation}
    \bold{B}_{rad}(\bold{r}, \theta, t) = - \frac{\mu_0 \ddot{p}(t_0)}{4 \pi c} \left(\frac{\sin \theta}{r}\right) \hat{\boldsymbol{\phi}}
\end{equation}
És a Poynting-vektor:
\begin{equation}
    \bold{S}_{rad}(\bold{r}, \theta, t) \approx \frac{1}{\mu_0} \bold{E}_{rad} \times \bold{B}_{rad} = \frac{\mu_0}{16 \pi^2 c} \left[\ddot{p}(t_0)\right]^2 \left(\frac{\sin^2 \theta}{r^2}\right) \hat{\boldsymbol{r}}
\end{equation}
A teljesítmény pedig, ami átmegy a gömbfelületen:
\begin{equation}
    P_{rad}(t) = \oint \bold{S}_{rad}(\bold{r}, \theta, t) \cdot d\bold{A} = \frac{\mu_0}{6 \pi c} \left[\ddot{p}(t - \frac{r}{c})\right]^2
\end{equation}
A teljes kisugárzott teljesítmény tehát csak a dipólusmomentum második időderiváltjától függ:
\begin{equation}
    P_{rad}(t) \approx \frac{\mu_0}{6 \pi c} \left[\ddot{p}(t - \frac{r}{c})\right]^2
\end{equation}

\subsubsection*{Példa: egy töltés sugárzása}
Ha egy töltés $q$ egyenes vonal mentén oszcillál, akkor a dipólusmomentuma:
\begin{equation}
    \bold{p}(t) = q \bold{r}(t)
\end{equation}
A második időderivált pedig:
\begin{equation}
    \ddot{\bold{p}}(t) = q \ddot{\bold{r}}(t) = q \bold{a}(t)
\end{equation}
Tehát a kisugárzott teljesítmény:
\begin{equation}
    P_{rad}(t) = \frac{\mu_0 q^2 a^2}{6 \pi c}
\end{equation}
Ez a híres \textbf{Larmor formula}, ami megadja a kisugárzott teljesítményt egy gyorsuló töltés esetében. Megjegyzendő, hogy a kisugárzott teljesítmény arányos a gyorsulás négyzetével, tehát egy állandó sebességgel mozgó töltés nem sugároz.
Ez összhangban van azzal a ténnyel, hogy egy egyenes vonalú, egyenletes mozgás nem generál időben változó elektromágneses teret.

\end{document}